% Major project report -- compile with:  xelatex report && bibtex report && xelatex report && xelatex report
\documentclass[12pt,oneside,openany]{report}

\usepackage{pdfpages}
\usepackage{setspace}
\usepackage{mathrsfs}
\usepackage{amsmath,amssymb,amsthm}
\usepackage{mathtools}
\usepackage[no-math]{mathspec}
\IfFontExistsTF{Times New Roman}{%
  \setmainfont{Times New Roman}
  \setmathsfont(Digits,Latin,Greek){Times New Roman}}{%
  \setmainfont{TeX Gyre Termes}
  \setmathsfont(Digits,Latin,Greek){TeX Gyre Termes}}

\usepackage[toc,page]{appendix}
\usepackage{float}
\graphicspath{{Figures/}}
\usepackage{wrapfig}
\usepackage{xcolor}

\newenvironment{chapterabstract}{%
  \par\nobreak\noindent
  \textbf{\textit{Summary}\hrulefill}\par\nobreak
  \small
  \noindent\ignorespaces
}{%
  \par\nobreak\normalsize
  \vskip-\ht\strutbox\noindent
  \textbf{\hrulefill}%
}

\usepackage[font=small]{caption}
\usepackage[a4paper,left=1in,right=1in,top=1in,bottom=1in,footskip=.25in]{geometry}
\usepackage{graphicx}
\usepackage{multirow}
\usepackage{booktabs}
\usepackage{enumitem}
\usepackage{tikz}
\usetikzlibrary{calc}
\usepackage{pgfplots}
\pgfplotsset{compat=1.17}
\usepgfplotslibrary{groupplots}
\usepackage{natbib}
\usepackage[english]{babel}
\usepackage{hyperref}
\hypersetup{
    colorlinks=true,
    linkcolor=blue,
    filecolor=magenta,
    urlcolor=cyan,
    citecolor=blue,
    pdftitle={Exact k-Dispersion for Points in Convex Position},
}
\usepackage[capitalise,noabbrev]{cleveref}

\setlength{\arrayrulewidth}{.3mm}

% ---------------------------------------------------------------------------
% Theorem environments and notation
\newtheorem{theorem}{Theorem}[chapter]
\newtheorem{lemma}[theorem]{Lemma}
\newtheorem{proposition}[theorem]{Proposition}
\newtheorem{corollary}[theorem]{Corollary}
\newtheorem{conjecture}[theorem]{Conjecture}
\theoremstyle{definition}
\newtheorem{definition}[theorem]{Definition}
\theoremstyle{remark}
\newtheorem{remark}[theorem]{Remark}

\floatstyle{ruled}
\newfloat{algorithm}{tbp}{loa}[chapter]
\floatname{algorithm}{Algorithm}
\crefname{algorithm}{Algorithm}{Algorithms}
\floatstyle{plain}
\restylefloat{figure}
\restylefloat{table}

\newcommand{\R}{\mathbb{R}}
\newcommand{\far}{\mathrm{far}}
\newcommand{\per}{\mathrm{per}}
\newcommand{\conv}{\mathrm{conv}}
\newcommand{\ang}{\angle}
\newcommand{\rstar}{r^{*}}
\newcommand{\circlethree}[3]{\bigcirc(#1,#2,#3)}
\newcommand{\Arc}{\mathrm{A}}
\newcommand{\Tr}{T_r}
\newcommand{\nil}{\bot}
\newcommand{\AlgCapped}{\textsc{CappedDecision}}
\newcommand{\AlgCert}{\textsc{Certified}}
\newcommand{\code}[1]{\nolinkurl{#1}}
\makeatletter\newcommand{\tabinput}[1]{\@@input #1.tex }\makeatother
\newcommand{\nVerify}{320}
\newcommand{\nVerifyOk}{320}
\newcommand{\nMedium}{64}
\newcommand{\nMediumOk}{64}
\newcommand{\nMediumBrute}{39}
\newcommand{\nVerifyPy}{320}
\newcommand{\nVerifyFmm}{320}
\newcommand{\nScaling}{462}
\newcommand{\nScalingOk}{462}
\newcommand{\expCappedFour}{3.29}
\newcommand{\expCappedEight}{3.45}
\newcommand{\expCappedSixteen}{3.52}
\newcommand{\expCappedThirtyTwo}{3.51}
\newcommand{\windowMedian}{6.0}
\newcommand{\windowMax}{34.7}
\newcommand{\windowSaved}{4.3}
\newcommand{\windowSpeedup}{1.35}
\newcommand{\farTripleMedian}{51}
\newcommand{\farTripleMin}{3}
\newcommand{\farTripleMax}{95}
\newcommand{\capSpeedupKFour}{1.8}
\newcommand{\tripleSpeedupHundred}{3.6}
\newcommand{\certSpeedup}{5.2}
\newcommand{\certSpeedupLo}{3.5}
\newcommand{\certSpeedupHi}{6.8}
\newcommand{\certOverhead}{11}
\newcommand{\certScalingNo}{39}
\newcommand{\certTotal}{846}
\newcommand{\certYes}{807}
\newcommand{\certTwochainNo}{30}
\newcommand{\certTwochainAll}{77}
\newcommand{\huntCone}{55{,}746{,}393}
\newcommand{\huntVfour}{55{,}746{,}393}

\newcommand{\arcSmallInst}{2{,}701}
\newcommand{\arcDchecks}{2{,}648{,}328}
\newcommand{\arcOptChecks}{23{,}927}
\newcommand{\arcBigInst}{176}
\newcommand{\arcBigChecks}{1{,}798}
\newcommand{\arcUncheckedWrong}{8}
\newcommand{\arcPyAtwo}{23{,}263}
\newcommand{\arcHoneFourFive}{31{,}532{,}602}
\newcommand{\arcHtwo}{2{,}350{,}993}
\newcommand{\arcSlope}{1.15}
\newcommand{\arcSlopeOone}{2.11}
\newcommand{\arcMillion}{6}
\newcommand{\arcSpeedupDP}{11943}


% ---------------------------------------------------------------------------
% Project details -- EDIT THESE
\newcommand{\ReportTitle}{Exact $k$-Dispersion for Points in Convex Position: Search Trees, Greedy Structure and a Capped Delaunay Dynamic Program}
\newcommand{\StudentA}{Akshat Bharara}
\newcommand{\StudentB}{Rakshith Ashok Kumar}
\newcommand{\StudentC}{Utsav Singh Bhamra}
\newcommand{\RegA}{231CS110}
\newcommand{\RegB}{231CS147}
\newcommand{\RegC}{231CS163}
\newcommand{\Guide}{Prof.\ Manjanna B}
\newcommand{\GuideDesignation}{Professor}
\newcommand{\Department}{Computer Science and Engineering}
\newcommand{\SubmissionDate}{November 2026}
% ---------------------------------------------------------------------------

\pagenumbering{roman}
\begin{document}
\emergencystretch=2em

\begin{titlepage}
\begin{center}
{\Large\bfseries EXACT $k$-DISPERSION FOR POINTS IN CONVEX POSITION:\\[4pt]
SEARCH TREES, GREEDY STRUCTURE AND A CAPPED\\[4pt] DELAUNAY DYNAMIC PROGRAM\par}

\vspace*{0.45in}
{\large Major Project Report}\\[6pt]
\textit{submitted in partial fulfilment of the requirements for the award of the degree of}\\[10pt]
{\large\bfseries BACHELOR OF TECHNOLOGY}\\[4pt]
in\\[4pt]
{\large\bfseries \Department}

\vspace*{0.4in}
\textit{by}\\[8pt]
\begin{tabular}{@{}l@{\hspace{2em}}l@{}}
\textbf{\StudentA} & \RegA\\
\textbf{\StudentB} & \RegB\\
\textbf{\StudentC} & \RegC
\end{tabular}

\vspace*{0.4in}
\textit{under the guidance of}\\[6pt]
{\bfseries \Guide}

\vfill
\IfFileExists{Figures/nitk_logo.png}{\includegraphics[width=1.3in]{nitk_logo.png}\\[12pt]}{}
{\large\bfseries DEPARTMENT OF \MakeUppercase{\Department}\\[4pt]
NATIONAL INSTITUTE OF TECHNOLOGY KARNATAKA\\[4pt]
SURATHKAL, MANGALORE -- 575025}\\[10pt]
{\large \SubmissionDate}
\end{center}
\end{titlepage}

\thispagestyle{empty}

\begin{spacing}{1.5}

%%%%%%%%%%%%%%%%%%%%%%%%%%%%%%%%%%%%%%%%%%%%%%%%%%%%%%%%%%%%%%%%%%%%
% Declaration
\newpage
\begin{titlepage}
\begin{center}
\large \textbf{DECLARATION} \\
\textit{By the B.Tech.\ Students}
\end{center}

\vspace*{0.4in}
\indent We hereby \textit{declare} that the Major Project Report entitled
{\bf\textquotedblleft\ReportTitle\textquotedblright}, which is being submitted
to the {\bf\textit{National Institute of Technology Karnataka, Surathkal}} in
partial fulfilment of the requirements for the award of the Degree of
{\bf\textit{Bachelor of Technology}} in {\bf\textit{\Department}} is a
\textit{bonafide report of the work carried out by us}. The material contained
in this report has not been submitted to any University or Institution for the
award of any degree.

\vspace*{0.5in}
\begin{singlespacing}
\noindent
\begin{tabular}{@{}l@{\hspace{3em}}l@{\hspace{3em}}l@{}}
\textbf{Name} & \textbf{Register No.} & \textbf{Signature}\\[6pt]
\StudentA & \RegA & \rule{1.3in}{0.4pt}\\[18pt]
\StudentB & \RegB & \rule{1.3in}{0.4pt}\\[18pt]
\StudentC & \RegC & \rule{1.3in}{0.4pt}
\end{tabular}

\vspace*{0.4in}
\noindent Department of \Department\\
National Institute of Technology Karnataka, Surathkal
\end{singlespacing}

\vspace*{0.5in}
\noindent Place: NITK, Surathkal\\
\noindent Date: \rule{1.2in}{0.4pt}
\end{titlepage}

%%%%%%%%%%%%%%%%%%%%%%%%%%%%%%%%%%%%%%%%%%%%%%%%%%%%%%%%%%%%%%%%%%%
% Certificate
\newpage
\begin{titlepage}
\begin{center}
\large \textbf{CERTIFICATE}
\end{center}

\vspace*{0.4in}
\indent This is to \textit{certify} that the Major Project Report entitled
{\bf\textquotedblleft\ReportTitle\textquotedblright}, submitted by
{\bf \StudentA} (\RegA), {\bf \StudentB} (\RegB) and {\bf \StudentC} (\RegC)
as the record of the work carried out by them, is \textit{accepted} as the
\textit{B.Tech.\ Major Project Report submission} in partial fulfilment of the
requirements for the award of the degree of {\bf\textit{Bachelor of
Technology}} in {\bf\textit{\Department}}.

\vspace*{0.9in}
\begin{singlespacing}
\hspace*{2.5in}
\parbox{3in}
{\noindent {\bf \Guide} \\
\noindent Project Guide \\
\noindent \GuideDesignation\\
\noindent Department of \Department\\
\noindent NITK Surathkal -- 575025
\vspace*{1.1in}\\
\noindent {\bf Head of the Department} \\
\noindent (Signature with Date and Seal)}
\end{singlespacing}
\end{titlepage}


%%%%%%%%%%%%%%%%%%%%%%%%%%%%%%%%%%%%%%%%%%%%%%%%%%%%%%%%%%%%%%%%%%%%%%%%%
% Abstract
\newpage
\begin{titlepage}
\pagestyle{empty}
\begin{center}
\large \textbf{ABSTRACT}
\end{center}

\begin{singlespacing}
Given $n$ points in convex position in the plane and an integer $k$, the
\emph{$k$-dispersion} problem asks for $k$ of the points whose minimum pairwise
distance is as large as possible. It models the placement of facilities that
must be kept apart, such as transmitters or hazardous plants along a coastline.
Singireddy, Basappa and Mitchell (Discrete Applied Mathematics, 2026) gave a
fixed-parameter algorithm that grows a solution from one point in both
directions along the hull; its correctness rests on a claim whose published
proof is informal. Tkachenko and Wang (SWAT 2026) gave an $O(n^3\log^2n)$
algorithm.

This project makes three contributions. First, we prove the claim of
Singireddy et al.\ for $k\le4$, which makes their algorithm correct for these
values of $k$; for every $k$ we show that the first step of the search can be
made safely, and we give point sets showing that natural rules for choosing the
starting point fail for $k=5,6,7$. Along the way we show that the structural
properties that make $k=3$ easy all fail for $k\ge4$. Second, we identify the
\emph{arc property}---for every point, the points far from it are contiguous
along the hull---under which a single greedy chain is exact; for unimodal point
sets this gives an $O(n\log^2n)$ expected-time algorithm that solves instances
with a million points in seconds. Third, a Delaunay certificate for
well-separated sets yields a simple exact algorithm that decides the problem in
$O(n^3\log k)$ time and solves it in $O(n^3\log k\log n)$ time.

All algorithms are implemented in C++ with exact integer arithmetic and were
validated against exhaustive search on thousands of instances.

\vspace*{12pt}
\noindent \textit{Keywords}: $k$-dispersion, obnoxious facility location,
convex position, Delaunay triangulation, dynamic programming, greedy
algorithms, fixed-parameter algorithms, computational geometry.
\end{singlespacing}
\end{titlepage}

%%%%%%%%%%%%%%%%%%%%%%%%%%%%%%%%%%%%%%%%%%%%%%%%%%%%%%%%%%%%%%%%%%%%%%%%%
% Acknowledgements
\newpage
\begin{titlepage}
\pagestyle{empty}
\begin{center}
\large \textbf{Acknowledgements}
\end{center}

We thank our project guide, \Guide, for guidance and encouragement throughout
this project. We thank the Head of the Department and the faculty of the
Department of \Department, NITK Surathkal, for providing the facilities and
support needed to carry out this work, and the members of the evaluation panel
for their comments. Finally, we thank our families and friends for their
constant support.

\vspace*{0.5in}
\hfill \StudentA\\
\hspace*{\fill} \StudentB\\
\hspace*{\fill} \StudentC
\end{titlepage}

%%%%%%%%%%%%%%%%%%%%%%%%%%%%%%%%%%%%%%%%%%%%%%%%%%%%%%%%%%%%%%%%%%%%
\newpage
\tableofcontents


%%%%%%%%%%%%%%%%%%%%%%%%%%%%%%%%%%%%%%%%%%%%%%%%%%%%%%%%%%%%%%%%%%%%
% Main matter
\clearpage
\pagenumbering{arabic}
\setcounter{page}{1}

\chapter{Introduction}\label{ch:intro}

\begin{chapterabstract}
This chapter introduces the $k$-dispersion problem, explains why points in
convex position are an interesting special case, states the problem
precisely, lists the objectives of the project and summarises its
contributions.
\end{chapterabstract}

\section{Background}

Facility location asks where to place a number of facilities so that some
objective is optimised. In the classical versions, such as $k$-center or
$k$-median, facilities should be \emph{close} to the clients they serve. In
\emph{obnoxious} facility location the opposite is true: facilities should be
far from one another, because they interfere with or endanger each other.
Typical examples are radio transmitters that must not interfere, franchise
outlets that should not compete for the same customers, and storage plants for
hazardous material that should be placed so that an accident at one plant does
not reach another~\citep{tamir,ravi}.

The most basic problem of this kind is \emph{max--min dispersion}, also called
$k$-dispersion or the maximally separated set problem~\citep{agarwal}: given a
set $P$ of $n$ candidate locations and an integer $k$, choose $k$ of them so
that the smallest distance between two chosen locations is as large as
possible. For points in the plane with Euclidean distances the problem is
NP-hard~\citep{wangkuo}; the best exact algorithms take $n^{O(\sqrt k)}$
time~\citep{akagi}, and a ratio of $1/2$ is the best known approximation
guarantee in polynomial time~\citep{ravi,tamir}.

A natural way to obtain efficient algorithms is to restrict the input. When the
candidate locations lie on the boundary of a convex region---for example, on
the coast of a convex island, the setting used to motivate the problem
in~\citet{sbm}---the points are in \emph{convex position}, and the problem
becomes solvable in polynomial time. This project studies $k$-dispersion for
points in convex position. Our starting point is the recent paper of
Singireddy, Basappa and Mitchell in \emph{Discrete Applied
Mathematics}~\citep{sbm}, which gives a fixed-parameter algorithm based on a
bounded search tree, a polynomial dynamic program and an approximation for
$k=3$.

\section{Problem statement}\label{sec:problem}

Let $P$ be a set of $n$ points in the plane in \emph{convex position}, that is,
every point of $P$ is a vertex of the convex hull $\conv(P)$. For an integer
$2\le k\le n$, the \emph{$k$-dispersion} problem asks for a subset
$S\subseteq P$ with $|S|=k$ that maximises
\[
  \delta(S)=\min_{p\ne q\in S}|pq| .
\]
We write $\rstar=\rstar_k(P)$ for the optimum. The \emph{decision version}
asks, for a given radius $r$, whether $P$ contains $k$ points whose pairwise
distances are all at least $r$ (an \emph{$r$-far $k$-set}). Since $\rstar$ is
one of the $\binom n2$ pairwise distances and the decision is monotone in $r$,
an algorithm for the decision version yields one for the optimisation version
by binary search.

\section{Motivation}

Two questions motivated the project.

\begin{enumerate}[leftmargin=*,itemsep=2pt]
\item \emph{Is the search tree of~\citet{sbm} correct?} Its correctness rests on
a claim (Claim~2 of~\citet{sbm}) that, if $k$ well-separated points exist, a
solution can be grown from some starting point by repeatedly adding the first
admissible point at one of the two ends of the chosen arc. The published proof
of this claim is a picture argument. A rigorous proof would put the
algorithm on solid ground; a counterexample would show that it is incorrect.
\item \emph{When is the problem easy?} For $k=3$ much faster algorithms are
known than for general $k$~\citep{kobayashi,tw25}. We wanted to understand
which structural properties make $k=3$ easy, whether they extend to larger
$k$, and for which inputs a simple greedy rule is exact.
\end{enumerate}

\section{Objectives of the project}

\begin{itemize}[leftmargin=*,itemsep=2pt]
\item Study the algorithms of~\citet{sbm} and the related work, and re-implement
them in exact integer arithmetic.
\item Prove or refute the claim on which the bounded search tree of~\citet{sbm}
relies, at least for small $k$.
\item Determine which greedy properties that hold for $k=3$ extend to larger
$k$, and identify a natural class of inputs on which greedy is exact.
\item Obtain a faster algorithm for such a special class, and a simple exact
algorithm for the general problem.
\item Implement all algorithms, validate them against exhaustive search, and
evaluate them experimentally.
\end{itemize}

\section{Contributions}\label{sec:contributions}

\begin{enumerate}[leftmargin=*,itemsep=2pt]
\item \textbf{The two-sided search tree} (\cref{sec:sbm}). We prove the claim
of~\citet{sbm} for $k\le4$ (\cref{thm:sbm4}), which makes their algorithm
correct for these $k$. For every $k$ we prove an attachment lemma that
certifies single steps of the search and show, using Delaunay ears, that the
first step can always be made safely (\cref{thm:firststep}). We also give point
sets showing that the starting point cannot be chosen by distances or angles
alone (\cref{prop:rootrules}).

\item \textbf{Greedy structure} (\cref{sec:greedy}). For $k=3$ every $r$-far
triple can be turned into one of a few canonical forms by a single
replacement, and an arc characterisation of $r$-far triples holds. For
$k\ge4$ these properties fail; in particular there are point sets in which no
solution contains a point together with the first point at distance at least
$r$ from it (\cref{thm:tightpair}). The $8$-point example for $k=4$ can be
checked by hand.

\item \textbf{The arc property} (\cref{sec:arc}). If, for every point $q$, the
points at distance at least $r$ from $q$ are contiguous along the hull, then a
single greedy chain decides the problem (\cref{thm:arc}). For unimodal point
sets, which include perturbed points on a circle, this gives an exact
algorithm with expected running time $O(n\log^2n)$ (\cref{thm:arcopt}).

\item \textbf{A simple exact algorithm} (\cref{ch:exact}). A set in convex
position is $r$-far if and only if it has a triangulation with long edges and
locally Delaunay diagonals (\cref{thm:cert-delaunay}). This certificate gives a
short correctness proof for a dynamic program over chords; capping the
per-chord information at $k-2$ thresholds decides the problem in
$O(n^3\log k)$ time and solves it in $O(n^3\log k\log n)$ time.

\item \textbf{Implementation and evaluation} (\cref{ch:results}). All
algorithms are implemented in C++ with exact integer predicates, cross-checked
against exhaustive search and independent code, and evaluated on instances
with up to $10^6$ points.
\end{enumerate}

\section{Organisation of the report}

\textbf{Chapter~\ref{ch:lit}} reviews the literature.
\textbf{Chapter~\ref{ch:prelim}} fixes notation and proves the Delaunay
certificate. \textbf{Chapter~\ref{sec:greedy}} studies greedy structure for
$k=3$ and $k\ge4$. \textbf{Chapter~\ref{sec:sbm}} analyses the two-sided search
tree of~\citet{sbm} and proves its correctness for $k\le4$.
\textbf{Chapter~\ref{sec:arc}} treats point sets with the arc property.
\textbf{Chapter~\ref{ch:exact}} presents the exact algorithm.
\textbf{Chapter~\ref{ch:results}} describes the implementation and the
experimental results, and \textbf{Chapter~\ref{ch:conclusion}} concludes with
limitations and future work.

\chapter{Literature Survey}\label{ch:lit}

\begin{chapterabstract}
This chapter surveys exact and approximation algorithms for dispersion,
with emphasis on points in convex position, and describes the two papers on
which the project builds: the bounded search tree and dynamic program of
Singireddy, Basappa and Mitchell, and the independent-set approach of
Tkachenko and Wang.
\end{chapterabstract}

\section{Dispersion in general}

For a set of $n$ points with arbitrary distances, $k$-dispersion is NP-hard
even when the distances satisfy the triangle inequality, and for Euclidean
points in the plane it is NP-hard as well~\citep{wangkuo}. Ravi, Rosenkrantz
and Tayi~\citep{ravi} showed that no polynomial algorithm achieves a ratio
better than $1/2$ for metric distances unless $\mathrm{P}=\mathrm{NP}$, and gave a
$1/2$-approximation; the same guarantee had been obtained earlier by
Tamir~\citep{tamir}. The algorithm is the farthest-first greedy rule, which is
also the classical $2$-approximation for $k$-center~\citep{gonzalez}.

For exact algorithms in the plane, Akagi et al.~\citep{akagi} gave an
$n^{O(\sqrt k)}$-time algorithm, together with linear-time algorithms for the
special cases in which the points lie on a line or on a circle; Araki and
Nakano~\citep{arakinakano} improved the line case to $O(\log n)$ time after
preprocessing. Chen and Dumitrescu~\citep{chendumitrescu} recently gave simple
algorithms for small~$k$. Agarwal, Overmars and Sharir~\citep{agarwal} studied
maximally separated sets in the plane, and Fekete and
Meijer~\citep{feketemeijer} considered dispersion with an average-distance
objective.

\section{Dispersion for points in convex position}

For $k=3$, Horiyama et al.~\citep{horiyama} gave an $O(n^2\log n)$ algorithm
for the Euclidean metric, and Kobayashi et al.~\citep{kobayashi} an $O(n^2)$
algorithm for points in convex position. Mishra, Rao and
Das~\citep{mishra} studied an approximation for points on a convex polygon.

\paragraph{Singireddy, Basappa and Mitchell.}
The paper of Singireddy, Basappa and Mitchell~\citep{sbm} is the starting point
of this project. It contains three results.
\begin{enumerate}[leftmargin=*,itemsep=2pt]
\item A \emph{bounded search tree}. The first point is placed at a vertex of the
hull; afterwards the chosen points always form a contiguous arc, and the next
point is the first vertex beyond one of the two ends of this arc that is far
enough from all chosen points. Branching on the end gives at most $2^{k}$
paths per starting vertex, and binary search over the candidate distances
gives an $O(2^kn^2\log^2n)$ algorithm, which is fixed-parameter tractable in
$k$. Its correctness rests on the claim that some starting vertex and some
sequence of ends reach $k$ points whenever a solution exists (Claim~2
of~\citep{sbm}). This claim is the subject of \cref{sec:sbm}.
\item A \emph{dynamic program} over Delaunay triangulations of the solution,
running in $O(n^4k^2)$ time. It was the first polynomial algorithm for
general~$k$.
\item An $O(\log n)$-time $1/(2\sqrt2)$-approximation for $k=3$ that uses only the
extreme points of the polygon.
\end{enumerate}

\paragraph{Tkachenko and Wang.}
Tkachenko and Wang~\citep{tw25} observed that the decision version of
dispersion is a maximum-independent-set problem in the unit-disk graph of the
points, and solved it by dynamic programming over triangulations; for
dispersion their algorithm runs in $O(n^{37/11}\log n)$ time, and for $k=3$
the decision problem can be solved in $O(n\log n)$ time. In a recent
paper~\citep{tw26} they compute a maximum-weight independent set of $n$ disks
in convex position in $O(n^3\log n)$ time, which gives an $O(n^3\log^2 n)$
algorithm for dispersion. A key ingredient~\citep[Lemma~12]{tw26} is that, for
a fixed chord, the constraints imposed by the triangle on the other side of
the chord are monotone along the bisector of the chord; all subproblems of a
chord can therefore be evaluated by sorting and binary search. The dynamic
program of \cref{ch:exact} uses the same monotonicity.

\section{Summary of known bounds}

\Cref{tab:known} summarises the exact algorithms for points in convex position,
together with the results of this project.

\begin{table}[htb]
\centering
\caption{Exact algorithms for $k$-dispersion of $n$ points in convex position.
The general bound of~\citet{tw25} is as reported in~\citet{tw26}.}
\label{tab:known}
\small\setlength{\tabcolsep}{4pt}
\begin{tabular}{@{}lll@{}}
\toprule
Reference & Time & Remarks\\
\midrule
\citet{sbm} & $O(2^k n^2\log^2 n)$ & search tree; correct for $k\le4$ by \cref{thm:sbm4}\\
\citet{sbm} & $O(n^4k^2)$ & Delaunay dynamic program\\
\citet{tw25} & $O(n^{37/11}\log n)$ & via unit-disk independent set\\
\citet{tw26} & $O(n^3\log^2 n)$ & via weighted independent set of disks\\
\citet{kobayashi} & $O(n^2)$ & $k=3$\\
\citet{tw25} & $O(n\log n)$ & $k=3$, decision version\\
\citet{akagi} & $O(n)$ after sorting & points on a circle\\
\midrule
\textbf{This project} (\cref{thm:arcopt}) & $O(n\log^2n)$ expected & unimodal point sets\\
\textbf{This project} (\cref{thm:opt}) & $O(n^3\log k\log n)$ & decision in $O(n^3\log k)$\\
\bottomrule
\end{tabular}
\end{table}

\section{Research gaps}

The survey leaves the following gaps, which this project addresses.
\begin{itemize}[leftmargin=*,itemsep=2pt]
\item The correctness of the fixed-parameter algorithm of~\citet{sbm} depends on
a claim without a complete proof.
\item The fast algorithms for $k=3$ exploit structure tied to the first far
points of a point; it was not known which of these properties survive for
larger~$k$.
\item Apart from points on a circle, no class of point sets in convex
position was known on which dispersion can be solved in near-linear time.
\item The cubic algorithms for general $k$ pay a $\log n$ factor per decision even
when $k$ is small.
\end{itemize}

\chapter{Preliminaries}\label{ch:prelim}
\section{Notation and conventions}\label{sec:prelim}

Throughout, $P=(p_0,\dots,p_{n-1})$ is a sequence of $n\ge3$ points in
\emph{strictly} convex position (no three collinear), listed counterclockwise
(ccw); indices are taken modulo~$n$. We write $|pq|$ for the Euclidean
distance. For $i\ne j$, $P(i,j)$ denotes the points strictly between $p_i$ and
$p_j$ when going ccw from $p_i$ to $p_j$. For $r>0$, a set is \emph{$r$-far}
if all its pairwise distances are at least $r$, and $\far_r(i,j)$ is the
predicate $|p_ip_j|\ge r$. For three points, $\ang abc\in(0,\pi)$ is the angle
at $b$ in the triangle $abc$, and $\circlethree abc$ is the circle through
them. $\per(P)$ is the perimeter of $\conv(P)$.

A \emph{triangulation} of a set $S$ in convex position, $|S|\ge3$, is a
maximal set of pairwise non-crossing segments (\emph{edges}) between points of
$S$; its edges are the sides of $\conv(S)$ and the \emph{diagonals}. A
diagonal $ab$ with incident triangles $abc$ and $bad$ is \emph{locally
Delaunay} if $d$ is not strictly inside $\circlethree abc$. Since $acbd$ is a
convex quadrilateral, this is equivalent to
\begin{equation}\label{eq:ld}
\ang acb+\ang adb\le\pi .
\end{equation}
A triangulation is \emph{Delaunay} if no point of $S$ lies strictly inside the
circumcircle of any of its triangles. Delaunay triangulations exist for every
finite point set not contained in a line (triangulate the cells of the
Delaunay subdivision arbitrarily).

We repeatedly use two elementary facts.
\begin{description}[leftmargin=1.5em,itemsep=1pt]
\item[(Q)] In a convex quadrilateral $abcd$ the diagonals are together longer
than either pair of opposite sides: $|ac|+|bd|>|ab|+|cd|$ and
$|ac|+|bd|>|bc|+|da|$. (Apply the triangle inequality at the crossing point of
the diagonals.)
\item[(A)] Seen from a vertex $v$ of a convex polygon, the other vertices
appear in their cyclic order within an angle smaller than $\pi$. Hence, if
$v,u_1,u_2,u_3$ occur in this ccw order, $\ang u_1vu_3=\ang u_1vu_2+\ang u_2vu_3$.
\end{description}

\paragraph{Model of computation and exact predicates.}
Our algorithms compare distances and evaluate the in-circle predicate; no other
geometric primitive is used. The analysis holds on a real RAM, where the term
$n^3/w$ below is to be read as $n^3$, and on a word RAM with word size
$w=\Omega(\log n)$ for integer input. For integer coordinates of absolute value
at most $2^{29}$, squared distances fit into $62$ bits and the in-circle
determinant is smaller than $2^{124}$ in absolute value, so that signed
$128$-bit integers suffice; this is what our implementation does. All angle
comparisons below reduce to in-circle tests by the inscribed angle theorem:
\begin{description}[leftmargin=1.5em,itemsep=1pt]
\item[(I1)] if $x$ and $y$ lie on the same side of the line $ab$, then
$\ang axb>\ang ayb$ iff $x$ is strictly inside $\circlethree ayb$;
\item[(I2)] if $x$ and $c$ lie on opposite sides of the line $ab$, then
$\ang axb+\ang acb>\pi$ iff $x$ is strictly inside $\circlethree acb$.
\end{description}

\section{A Delaunay certificate for \texorpdfstring{$r$}{r}-far sets}\label{sec:certificate}

\begin{lemma}[Delaunay lemma {\citep{delaunay,edelsbrunner}}]\label{lem:delaunay}
If every diagonal of a triangulation $T$ of a finite point set $S$ is locally
Delaunay, then $T$ is a Delaunay triangulation of $S$.
\end{lemma}

It is well known that a closest pair of a point set is joined by an edge of
its Delaunay triangulation. We need the statement for an \emph{arbitrary}
Delaunay triangulation in the presence of cocircular points, and give the
short proof for sets in convex position.

\begin{lemma}\label{lem:closest}
Let $T$ be a Delaunay triangulation of a set $S$ in strictly convex position,
$|S|\ge3$. Then every closest pair of $S$ is joined by an edge of $T$.
\end{lemma}
\begin{proof}
Let $p,q$ be a closest pair and suppose that $pq$ is not an edge of $T$. Then
$pq$ is not a side of $\conv(S)$, and it does not pass through a third point
of $S$. Hence, starting at $p$, the segment $pq$ runs through the interior of
a triangle $puv$ of $T$ and leaves it through an interior point of the side
$uv$. So $u$ and $v$ lie strictly on opposite sides of the line $pq$, and
$p,u,q,v$ are the vertices of a convex quadrilateral in this cyclic order,
with diagonals $pq$ and $uv$. Because $T$ is Delaunay, $q$ is not strictly
inside $\circlethree puv$, that is, $\ang upv+\ang uqv\le\pi$ by~(I2). The
angles of the quadrilateral sum to $2\pi$, so $\ang puq+\ang pvq\ge\pi$, and
one of them, say $\ang puq$, is at least $\pi/2$. Then $pq$ is the strictly
longest side of the triangle $puq$, so $|pu|<|pq|$, contradicting the choice
of $p,q$.
\end{proof}

\begin{theorem}[Delaunay certificate]\label{thm:cert-delaunay}
Let $S$ be a set in strictly convex position, $|S|\ge3$, and let $r>0$. Then
$S$ is $r$-far if and only if $S$ has a triangulation $T$ such that
\begin{enumerate}[label=(\alph*),itemsep=0pt,topsep=2pt]
\item every edge of $T$ has length at least $r$, and
\item every diagonal of $T$ is locally Delaunay.
\end{enumerate}
\end{theorem}
\begin{proof}
($\Leftarrow$) By (b) and \cref{lem:delaunay}, $T$ is a Delaunay
triangulation. By \cref{lem:closest} the minimum pairwise distance of $S$ is
the length of an edge of $T$, which is at least $r$ by~(a).
($\Rightarrow$) Take any Delaunay triangulation $T$ of $S$. It satisfies (b)
by definition, and (a) because all its edges join points of the $r$-far
set~$S$.
\end{proof}

Conditions (a) and (b) are \emph{local}: (a) concerns single edges and (b)
pairs of adjacent triangles. The dynamic program of
Tkachenko and Wang~\citep{tw26} is built on the same idea of gluing triangles
under a local compatibility condition (for disks it is phrased with
Voronoi vertices of triples); \cref{thm:cert-delaunay} isolates the
statement that makes such a program correct and reduces its proof to the
Delaunay lemma.


\chapter{Greedy structure: true for \texorpdfstring{$k=3$}{k=3}, false for \texorpdfstring{$k\ge4$}{k>=4}}\label{sec:greedy}

For $k=3$ the decision problem can be solved in near-linear
time~\citep{tw25}, and $O(n^2)$ time suffices for the optimisation
problem~\citep{kobayashi}. These algorithms exploit that an $r$-far triple, if
one exists, can be found next to the ``first far points'' of some point. It is
natural to ask how much of this structure survives for larger $k$, because any
property of this kind immediately yields a fast greedy or bounded-search
algorithm. This section gives a complete answer for the properties we are
aware of.

Fix $r>0$. For $q\in P$ let $a(q)$ be the first point after $q$ in ccw order
with $|q\,a(q)|\ge r$, and $b(q)$ the first such point in clockwise order;
both are undefined if no point is at distance at least $r$ from $q$. Note that
the distance from $q$ is in general not unimodal along the hull, so points
after $a(q)$ may again be close to $q$. Let $\Arc(q)$ be the set of points on
the ccw arc from $a(q)$ to $b(q)$ (which does not contain $q$). We call
$\{q,a(q)\}$ a \emph{tight pair}.

\begin{lemma}\label{lem:arcbasic}
Let $S$ be an $r$-far set with $|S|\ge2$, and let $s,s'$ be consecutive points
of $S$ in ccw order. Then $a(s)$ lies in $P(s,s')\cup\{s'\}$, and
$S\setminus\{s\}\subseteq\Arc(s)$.
\end{lemma}
\begin{proof}
All points of $S\setminus\{s\}$ are at distance at least $r$ from $s$; the first
of them in ccw order is $s'$ and the first in clockwise order is the
predecessor of $s$ in $S$.
\end{proof}

We consider the following properties of $(P,k,r)$.
\begin{description}[leftmargin=3.2em,labelwidth=2.6em,itemsep=2pt]
\item[$(\mathrm{G}_k)$] \emph{Arc characterisation:} $P$ has an $r$-far $k$-set
if and only if $\Arc(q)$ contains an $r$-far $(k-1)$-set for some $q\in P$.
\item[$(\mathrm{TT}_k)$] \emph{Tight triple:} if $P$ has an $r$-far $k$-set,
then it has one that contains $q$, $a(q)$ and $b(q)$ for some $q$.
\item[$(\mathrm{TP}_k)$] \emph{Tight pair:} if $P$ has an $r$-far $k$-set, then
it has one that contains a tight pair $\{q,a(q)\}$.
\item[$(\mathrm{ST}_k)$] \emph{Seeded search tree:} if $P$ has an $r$-far $k$-set, then
for some $q$ the following search reaches $k$ points: start with
$\{q,a(q)\}$, a ccw frontier at $a(q)$ and a clockwise frontier at $q$, and
repeatedly add, on a side of the search's choice, the first point beyond that
side's frontier (and before the other frontier) that is at distance at least
$r$ from all chosen points; the added point becomes the new frontier.
\item[$(\mathrm{AG}_k)$] \emph{Alternating greedy:} as $(\mathrm{ST}_k)$, but
the sides strictly alternate (either side may start).
\end{description}
Every set produced by a search of $(\mathrm{ST}_k)$ is $r$-far and contains
$\{q,a(q)\}$, so
$(\mathrm{AG}_k)\Rightarrow(\mathrm{ST}_k)\Rightarrow(\mathrm{TP}_k)$, and
trivially $(\mathrm{TT}_k)\Rightarrow(\mathrm{TP}_k)$. If $(\mathrm{ST}_k)$ held
for all point sets, the decision problem could be solved in $O(2^kn^2k)$ time
by exploring, for each of the $n$ seeds, a binary tree of depth $k-2$ whose
nodes cost $O(nk)$; $(\mathrm{AG}_k)$ would even give $O(n^2k)$. Property
$(\mathrm{TP}_k)$ is the weakest statement that ties a solution to the first
far points, and therefore the most robust candidate. The bounded search tree
of Singireddy, Basappa and Mitchell~\citep{sbm} is a two-sided variant of
$(\mathrm{ST}_k)$ that is seeded with a single point; we study it separately
in \cref{sec:sbm}.

\section{Triples}\label{sec:k3}

\begin{theorem}\label{thm:k3}
Let $\{x,y,z\}$ be an $r$-far triple of $P$, in ccw order.
\begin{enumerate}[label=(\roman*),itemsep=1pt,topsep=2pt]
\item At least one of $\{x,a(x),b(x)\}$, $\{y,a(y),b(y)\}$,
$\{z,a(z),b(z)\}$ is an $r$-far triple.
\item At least one of $\{x,a(x),z\}$, $\{y,a(y),x\}$, $\{z,a(z),y\}$ is an
$r$-far triple.
\end{enumerate}
In particular $(\mathrm{TT}_3)$, $(\mathrm{TP}_3)$, $(\mathrm{ST}_3)$ and
$(\mathrm{AG}_3)$ hold for every point set in convex position and every $r$.
\end{theorem}
\begin{proof}
(i) By \cref{lem:arcbasic}, the ccw order around the hull is
$x,a(x),y,z,b(x)$, where $a(x)=y$ and $b(x)=z$ are allowed, and
$a(x)\ne b(x)$. Suppose that $|a(x)b(x)|<r$. In the triangle $x\,a(x)\,b(x)$
the side $a(x)b(x)$ is then strictly shorter than the other two, so the
opposite angle is smaller than $\pi/3$, and by (A),
$\ang yxz\le\ang a(x)\,x\,b(x)<\pi/3$. If the triples of $y$ and $z$ also
failed to be $r$-far, the same argument would give $\ang zyx<\pi/3$ and
$\ang xzy<\pi/3$, and the angles of the triangle $xyz$ would sum to less than
$\pi$.

(ii) If $a(x)=y$ the first triple is $\{x,y,z\}$. Otherwise $a(x)\in P(x,y)$
and $|x\,a(x)|\ge r$, so the first triple is $r$-far unless $|a(x)z|<r$. In
that case $a(x)z$ is the strictly shortest side of the triangle $x\,a(x)\,z$,
hence $\ang a(x)\,x\,z<\pi/3$, and by (A) applied to $x,a(x),y,z$ we get
$\ang yxz<\pi/3$. Likewise, failure of the second triple gives
$\ang zyx<\pi/3$ and failure of the third gives $\ang xzy<\pi/3$. All three
cannot fail.

The tight triple of (i) contains the tight pair $\{q,a(q)\}$. For
$(\mathrm{ST}_3)$ and $(\mathrm{AG}_3)$, take the triple $\{q,a(q),z'\}$
provided by (ii): the search started from $\{q,a(q)\}$ needs a single step,
and its ccw step succeeds because the arc beyond $a(q)$ contains a point,
namely $z'$, at distance at least $r$ from both seeds.
\end{proof}

Part~(i) says that $P$ has an $r$-far triple if and only if
$|a(q)b(q)|\ge r$ for some $q$, a characterisation in the spirit of the
fast algorithms for $k=3$~\citep{kobayashi,tw25}; we include the proof
because it is short and because part~(ii) uses the same idea. The arc
characterisation holds for triples as well.

\begin{proposition}\label{prop:g3}
The implication ``$\Rightarrow$'' of $(\mathrm{G}_k)$ holds for every
$k\ge2$, and $(\mathrm{G}_3)$ holds.
\end{proposition}
\begin{proof}
If $S$ is an $r$-far $k$-set and $q\in S$, then $S\setminus\{q\}$ is an
$r$-far $(k-1)$-set contained in $\Arc(q)$ by \cref{lem:arcbasic}.

For $k=3$, let $x,y\in\Arc(q)$ with $|xy|\ge r$, write $a=a(q)$, $b=b(q)$,
and let the ccw order be $q,a,x,y,b$ with $a=x$ or $y=b$ allowed. Suppose that
$P$ has no $r$-far triple. Since $a\ne b$ and $|qa|,|qb|\ge r$, we have
$|ab|<r$. We distinguish cases according to the distances from $q$.
If $|qx|\ge r$ and $|qy|\ge r$, then $\{q,x,y\}$ is an $r$-far triple.
If $|qx|<r\le|qy|$, then $x\ne a$, and $|ay|<r$ because $\{q,a,y\}$ is not
$r$-far. In the convex quadrilateral $q,a,x,y$ the diagonals satisfy
$|qx|+|ay|<2r\le|qa|+|xy|$, contradicting~(Q). The case $|qy|<r\le|qx|$ is
symmetric: $y\neq b$, $|xb|<r$, and the quadrilateral $q,x,y,b$ has diagonals
$|qy|+|xb|<2r\le|xy|+|bq|$. Finally, if $|qx|<r$ and $|qy|<r$, then, as in the
proof of \cref{thm:k3}, $\ang xqy\le\ang aqb<\pi/3$, so $xy$ is not a longest
side of the triangle $qxy$ and $|xy|<\max\{|qx|,|qy|\}<r$, a contradiction.
\end{proof}

\section{Four or more points}\label{sec:k4}

\begin{proposition}\label{prop:gk}
$(\mathrm{G}_k)$ fails for every $k\ge4$: there is a set of $k$ points in
convex position and a radius $r$ such that no $r$-far $k$-set exists although
$\Arc(u)$ contains an $r$-far $(k-1)$-set for some point $u$.
\end{proposition}
\begin{proof}
Let $k\ge4$, choose $k-2$ distinct non-zero integers $j$ with $|j|\le k$, at
least one positive and one negative, and let $H=3(k+1)$. Let
\[
u=(-H^2,0),\qquad v=(H^2,0),\qquad w_j=\bigl(H^2-9j^2,\ 3jH^2\bigr),
\]
and $r=3H^2$. All points lie on the boundary of the strictly convex region
$\{(x,y):|x|\le H^2-y^2/H^4\}$ and are therefore in strictly convex position;
$u$ is alone on the left boundary arc, and the right arc carries $v$ and the
$w_j$, with points above and below $v$. We have $|uv|=2H^2<r$, while any two
other points differ by at least $3H^2=r$ in their $y$-coordinates. So the
point set itself, the only $k$-subset, is not $r$-far. The two hull neighbours
of $u$ are the lowest and the highest $w_j$; they are at distance at least
$r$ from $u$, so they are $a(u)$ and $b(u)$, and $\Arc(u)$ consists of all
points except $u$, which form an $r$-far $(k-1)$-set.
\end{proof}

\noindent For $k=4$ the idea is visible in the rhombus $(-1,0)$, $(0,-3)$,
$(1,0)$, $(0,3)$ with $r^2=10$: the short diagonal is the only close pair.

\begin{figure}[tb]
\centering
\begin{tikzpicture}[scale=0.047]
  \coordinate (p0) at (16,48);   \coordinate (p1) at (-83,26);
  \coordinate (p2) at (-114,18); \coordinate (p3) at (-30,-39);
  \coordinate (p4) at (-16,-48); \coordinate (p5) at (83,-26);
  \coordinate (p6) at (114,-18); \coordinate (p7) at (30,39);
  \draw[gray!60] (p0)--(p1)--(p2)--(p3)--(p4)--(p5)--(p6)--(p7)--cycle;
  \draw[blue!70!black,thick] (p0)--(p2)--(p4)--(p6)--cycle;
  \draw[blue!70!black,thick] (p0)--(p4);
  \draw[blue!70!black] (p2)--(p6);
  \draw[red!75!black,densely dashed,thick] (p1)--(p4) (p5)--(p0) (p3)--(p0) (p7)--(p4);
  \foreach \p/\pos/\lab in {p0/above/{$p_0=(16,48)$}, p2/left/{$p_2=(-114,18)$}, p4/below/{$p_4=(-16,-48)$}, p6/right/{$p_6=(114,-18)$}}
    {\fill[blue!70!black] (\p) circle (38pt); \node[\pos] at (\p) {\small \lab};}
  \foreach \p/\pos/\lab in {p1/above left/{$p_1=(-83,26)$}, p3/below left/{$p_3=(-30,-39)$}, p5/below right/{$p_5=(83,-26)$}, p7/above right/{$p_7=(30,39)$}}
    {\draw[fill=white] (\p) circle (38pt); \node[\pos] at (\p) {\small \lab};}
\end{tikzpicture}
\caption{The point set $P_\star$ of \cref{thm:tightpair} (to scale),
$r^2=10240$. Solid blue: the only $r$-far $4$-set $\{p_0,p_2,p_4,p_6\}$ and
its six distances; $|p_0p_4|=r$. Each white point is the first far point
$a(\cdot)$ of the blue point before it and is closer than $r$ to another blue
point (dashed red), so it cannot replace anything.}
\label{fig:pstar}
\end{figure}

\begin{theorem}\label{thm:tightpair}
$(\mathrm{TP}_4)$ is false, even at the optimal radius. Let $P_\star$ be the
eight points
\begin{align*}
&p_0=(16,48), &&p_1=(-83,26), &&p_2=(-114,18), &&p_3=(-30,-39),\\
&p_4=(-16,-48), &&p_5=(83,-26), &&p_6=(114,-18), &&p_7=(30,39),
\end{align*}
which are in strictly convex position in this ccw order, and let
$r^2=10240$. Then $\rstar_4(P_\star)=r$, the set $\{p_0,p_2,p_4,p_6\}$ is the
only $r$-far $4$-set, and it contains no tight pair. Consequently
$(\mathrm{TT}_4)$, $(\mathrm{ST}_4)$ and $(\mathrm{AG}_4)$ are false as well.
Moreover, $(\mathrm{TP}_k)$ is false for $k=5$, $6$ and $7$; point sets with
$2k$ points are listed in \cref{app:instances}.
\end{theorem}
\begin{proof}
$P_\star$ is symmetric under $p\mapsto-p$, which maps $p_i$ to $p_{i+4}$. The
squared distances from $p_0,\dots,p_3$ are
\begin{center}\small
\begin{tabular}{@{}c|cccccccc@{}}
 & $p_0$ & $p_1$ & $p_2$ & $p_3$ & $p_4$ & $p_5$ & $p_6$ & $p_7$\\\midrule
$p_0$ & -- & \textbf{10285} & \textbf{17800} & 9685 & \textbf{10240} & 9965 & \textbf{13960} & 277\\
$p_1$ &  & -- & 1025 & 7034 & 9965 & \textbf{30260} & \textbf{40745} & \textbf{12938}\\
$p_2$ &  &  & -- & \textbf{10305} & \textbf{13960} & \textbf{40745} & \textbf{53280} & \textbf{21177}\\
$p_3$ &  &  &  & -- & 277 & \textbf{12938} & \textbf{21177} & 9684\\
\end{tabular}
\end{center}
and the remaining ones follow by symmetry; entries of at least $r^2$ are in
bold. The points at distance at least $r$ from $p_1$ are $p_0,p_5,p_6,p_7$,
and the only pairs among these at distance at least $r$ are $\{p_0,p_6\}$ and
$\{p_6,p_7\}$, so $p_1$ lies in no $r$-far $4$-set; by symmetry neither does
$p_5$. The points at distance at least $r$ from $p_3$ are $p_2,p_5,p_6$, and
$|p_5p_6|<r$, so $p_3$ and, by symmetry, $p_7$ lie in no $r$-far $4$-set
either. Hence the only candidate is $\{p_0,p_2,p_4,p_6\}$, which is $r$-far
with minimum distance $|p_0p_4|=r$; in particular $\rstar_4(P_\star)=r$. From
the table, $a(p_0)=p_1$, $a(p_2)=p_3$, $a(p_4)=p_5$ and $a(p_6)=p_7$. None of
these belongs to the set, so it contains no tight pair.

For $k=5,6,7$ the instances of \cref{app:instances} were verified by
exhaustive enumeration of all $k$-subsets in exact integer arithmetic.
\end{proof}

\paragraph{Anatomy of the counterexample.}
The mechanism is easy to describe (\cref{fig:pstar}). The quadrilateral
$p_0p_2p_4p_6$ is a parallelogram whose short diagonal $p_0p_4$ has length
exactly $r$ and makes an angle of less than $\pi/3$ with the sides $p_0p_2$
and $p_4p_6$. The first far point $a(p_0)=p_1$ sits just outside the side
$p_0p_2$ at distance slightly more than $r$ from $p_0$; the triangle
$p_0p_1p_4$ has two sides of length about $r$ enclosing an angle below
$\pi/3$, so $|p_1p_4|<r$ and $p_1$ is blocked by $p_4$. The point
$a(p_2)=p_3$ sits just before $p_4$ and is slightly closer to $p_0$ than
$p_4$ is, so it is blocked by $p_0$.

More generally, suppose $(\mathrm{TP}_k)$ fails for $(P,r)$ and let
$S=\{s_1,\dots,s_k\}$ be an $r$-far $k$-set in ccw order. By
\cref{lem:arcbasic}, $a_i=a(s_i)$ lies strictly between $s_i$ and $s_{i+1}$,
and since $(S\setminus\{s_{i+1}\})\cup\{a_i\}$ would contain the tight pair
$\{s_i,a_i\}$, each $a_i$ must be closer than $r$ to some $s_j$ with
$j\notin\{i,i+1\}$. So every counterexample contains a convex $2k$-gon
$s_1a_1s_2a_2\cdots s_ka_k$ with $\{s_1,\dots,s_k\}$ $r$-far, $|s_ia_i|\ge r$,
and every $a_i$ blocked in this way; in particular it has at least $2k$
points, and our instances are smallest possible. For $k=3$ no such hexagon
exists by \cref{thm:k3}(ii).



\chapter{The two-sided search tree}\label{sec:sbm}

The fixed-parameter algorithm of Singireddy, Basappa and Mitchell~\citep{sbm}
decides whether an $r$-far $k$-set exists by the following search. A
\emph{state} is a set $C$ of chosen points that occupies a contiguous arc of
the hull, together with the complementary open arc $F$ (the \emph{free arc}).
Initially $C=\{q\}$ for a \emph{root} $q$, and $F=P\setminus\{q\}$. A
\emph{ccw step} adds the first point of $F$, in ccw order from the ccw end of
$C$, that is at distance at least $r$ from every point of $C$; a \emph{cw
step} is the mirror image. The search branches on the side, so that each root
spawns a binary tree of depth $k-1$; it succeeds if some root and some
sequence of sides reaches $|C|=k$. Every set it reaches is $r$-far, so the
search is sound. Its completeness is the following claim (Claim~2
of~\citet{sbm}, rephrased).

\begin{description}[leftmargin=3.2em,labelwidth=2.6em,itemsep=2pt]
\item[$(\mathrm{SBM}_k)$] If $P$ has an $r$-far $k$-set, the two-sided search
succeeds from some root.
\end{description}

\noindent By \cref{thm:tightpair}, both sides are needed: on $P_\star$ every
search that starts with a ccw step fails, whereas the two-sided search
succeeds. In this section we prove $(\mathrm{SBM}_k)$ for $k\le4$, prove that
the first step can always be made safely for every $k$, and show that the
starting point cannot be chosen by distances or angles alone.

\paragraph{Witnesses.}
Fix an $r$-far $k$-set $T$. We follow the search together with a
\emph{witness} $W$: an $r$-far $k$-set with $C\subseteq W$ and $W\setminus
C\subseteq F$. Write $W\setminus C=\{t_1,\dots,t_m\}$ in ccw order from the
ccw end of $C$. Let $x$ be the point that a ccw step would add. Since $t_1$ is
at distance at least $r$ from $C$, $x$ is $t_1$ or lies before $t_1$. We call
the ccw step \emph{safe} if $(W\setminus\{t_1\})\cup\{x\}$ is $r$-far; then it
is the new witness, and the state after the step again has a witness.
Symmetrically for cw steps with $t_m$.

\begin{lemma}[attachment]\label{lem:attach}
If some $c\in C$ satisfies $|c\,t_1|\le|t_1t_j|$ for all $2\le j\le m$, then
the ccw step is safe. Symmetrically, if some $c\in C$ satisfies
$|c\,t_m|\le|t_mt_i|$ for all $1\le i\le m-1$, then the cw step is safe.
\end{lemma}
\begin{proof}
If $x=t_1$ there is nothing to show. Otherwise $x$ lies strictly between the
ccw end of $C$ and $t_1$, and it is at distance at least $r$ from all of $C$.
Suppose that $|x\,t_j|<r$ for some $2\le j\le m$. The points $c,x,t_1,t_j$ are
in convex position in this ccw order, so by (Q),
$|c\,t_1|+|x\,t_j|>|c\,x|+|t_1t_j|$, that is,
$|c\,t_1|>|t_1t_j|+(|c\,x|-|x\,t_j|)>|t_1t_j|$, a contradiction. Hence $x$ is
at distance at least $r$ from $t_2,\dots,t_m$ and from $C$, and
$(W\setminus\{t_1\})\cup\{x\}$ is $r$-far.
\end{proof}

\Cref{lem:attach} says that a step is safe whenever the next point of the
witness has a chosen point among its nearest neighbours in the witness. At the
root this can always be arranged.

\begin{lemma}\label{lem:nnadjacent}
Let $\mathcal T$ be a Delaunay triangulation of a set $S$ in strictly convex
position. Every $p\in S$ has a nearest neighbour in $S$ that is adjacent to $p$
in $\mathcal T$. Consequently, if $|S|\ge3$, some $p\in S$ has a nearest
neighbour among its two neighbours on the hull of $S$.
\end{lemma}
\begin{proof}
Let $q$ be a nearest neighbour of $p$, and suppose that no nearest neighbour
of $p$ is adjacent to it. As in the proof of \cref{lem:closest}, the segment
$pq$ crosses an edge $uv$ of a triangle $puv$ of $\mathcal T$, $p,u,q,v$ form a
convex quadrilateral, $\ang puq+\ang pvq\ge\pi$ because $\mathcal T$ is
Delaunay, and say $\ang puq\ge\pi/2$; then $|pu|<|pq|$, a contradiction. For
the second claim, a triangulation of a convex polygon with at least four
vertices has an \emph{ear}, a vertex of degree two whose neighbours in the
triangulation are its two hull neighbours; for $|S|=3$ every vertex is such a
vertex.
\end{proof}

\begin{theorem}[safe first step]\label{thm:firststep}
Let $T$ be an $r$-far $k$-set, $k\ge2$. There are adjacent points $e,\rho$ of
$T$ (adjacent on the hull of $T$) such that the first step of the search from
root $\rho$ towards $e$ is safe with respect to the witness $T$.
\end{theorem}
\begin{proof}
For $k=2$ every step is safe. Otherwise, by \cref{lem:nnadjacent} some $e\in T$
has a hull neighbour $\rho$ in $T$ as a nearest neighbour in $T$. Start at
$\rho$ and step towards $e$; then $e$ is $t_1$ (or $t_m$) and $c=\rho$
satisfies the hypothesis of \cref{lem:attach}.
\end{proof}

For $k=3$ a safe first step leaves a single point to add, and the next step
is trivially safe, so $(\mathrm{SBM}_3)$ follows; this is a short proof of a
special case of \cref{thm:k3}.

\section{Four points}
For $k=4$ the claim holds. The proof needs two geometric lemmas.

\begin{lemma}\label{lem:opposite}
Let $p,q,s,t$ be in convex position in this ccw order. If
$|qs|\ge\max\{|ps|,|qt|\}$, then the interior angles of the quadrilateral at
$p$ and $t$ sum to at least $\pi$.
\end{lemma}
\begin{proof}
In the triangle $pqs$, $|qs|\ge|ps|$ gives $\ang qps\ge\ang pqs$; in the
triangle $qst$, $|qs|\ge|qt|$ gives $\ang qts\ge\ang qst$. The interior angle
at $p$ is at least $\ang qps$ and the one at $t$ at least $\ang qts$, while
the interior angles at $q$ and $s$ are $\ang pqs$ and $\ang qst$. Hence
$\ang p+\ang t\ge\ang q+\ang s=2\pi-\ang p-\ang t$.
\end{proof}

\begin{lemma}\label{lem:core}
There are no points $a,X,Y,c,d$ in strictly convex position in this ccw order
with $|aX|,|ad|,|cd|\ge r$ and $|aY|,|Xd|<r$.
\end{lemma}
\begin{proof}
The chords $aY$ and $Xd$ cross, so by (Q) applied to $a,X,Y,d$,
$|Yd|<|aY|+|Xd|-|aX|<r$. Write $\alpha_1=\ang cYd$, $\alpha_2=\ang dYa$,
$\alpha_3=\ang aYX$, $\varepsilon_1=\ang Yad$, $\varepsilon_2=\ang XaY$,
$\eta=\ang adX$, $\mu=\ang XdY$, $\delta_1=\ang Ydc$, $\gamma=\ang Ycd$,
$\lambda=\ang YXd$, $\zeta=\ang dXa$ and $\xi=\ang aXY$. By (A),
$\ang aXY=\lambda+\zeta$, $\ang adY=\eta+\mu$, $\ang XYd=\alpha_2+\alpha_3$
and $\ang Xad=\varepsilon_1+\varepsilon_2$. We use: convexity at $Y$,
$\alpha_1+\alpha_2+\alpha_3<\pi$; the angle sums of the triangles $XYd$,
$aXd$, $aYd$, $Ycd$, $acd$; and, comparing opposite sides,
$\xi<\alpha_3$ (triangle $aXY$: $|aY|<|aX|$),
$\gamma<\alpha_1$ (triangle $Ycd$: $|Yd|<|cd|$),
$\eta+\mu<\alpha_2$ (triangle $aYd$: $|aY|<|ad|$), and
$\varepsilon_1+\varepsilon_2<\min\{\eta,\zeta\}$ (triangle $aXd$:
$|Xd|<|aX|,|ad|$), which gives $\varepsilon_1+\varepsilon_2<\pi/3$.

From the triangles $XYd$ and $aXd$,
$\xi+\eta+\mu=(\lambda+\mu)+(\zeta+\eta)
=(\pi-\alpha_2-\alpha_3)+(\pi-\varepsilon_1-\varepsilon_2)$, hence with
$\xi<\alpha_3$,
\[\eta+\mu>2\pi-\alpha_2-2\alpha_3-\varepsilon_1-\varepsilon_2.\]
In the triangle $acd$ the angle at $d$ is $\eta+\mu+\delta_1<\pi$, and in
$Ycd$, $\delta_1=\pi-\alpha_1-\gamma>\pi-2\alpha_1$; so $\eta+\mu<2\alpha_1<
2(\pi-\alpha_2-\alpha_3)$. Together, $\alpha_2<\varepsilon_1+\varepsilon_2$.
Finally, in the triangle $aYd$,
$\pi-\varepsilon_1=\alpha_2+(\eta+\mu)<2\alpha_2<2\varepsilon_1+2\varepsilon_2$,
so $\pi<3(\varepsilon_1+\varepsilon_2)<\pi$, a contradiction.
\end{proof}

\begin{theorem}\label{thm:sbm4}
$(\mathrm{SBM}_4)$ holds: if $P$ has an $r$-far $4$-set $T$, the two-sided
search succeeds from a vertex of $T$, and even a witness-preserving search does.
\end{theorem}
\begin{proof}
Let $T=(a,b,c,d)$ in ccw order. Call $(u,v)$ a \emph{good pair} if $u,v$ are
hull neighbours in $T$ and $u$ is a nearest neighbour of $v$ in $T$. Start at
$u$ and step towards $v$; by \cref{lem:attach} the step is safe. Let $w$ be
the other neighbour of $u$ and $z$ the other neighbour of $v$; they are the
two remaining witness points, and $wz$ is an edge of $T$. By \cref{lem:attach},
applied with the chosen point $u$, one of the two next steps is safe unless $|wz|<|uw|$ and
$|wz|<|uz|$; after a safe second step one witness point remains, and the last
step always succeeds. So it suffices to find a good pair whose opposite edge
$wz$ is not the strictly shortest side of the triangle $uwz$.

Let a closest pair of $T$ be given. If it is a diagonal, say $ac$, then $ac$
is the diagonal of every Delaunay triangulation of $T$ (\cref{lem:closest}),
so $b$ is an ear and by \cref{lem:nnadjacent} a nearest neighbour of $b$ is
$a$ or $c$. For the good pair $(a,b)$ the triangle is $adc$ and $|dc|\ge|ac|$;
for $(c,b)$ it is $cda$ and $|da|\ge|ca|$. Either way we are done. If the
closest pair is an edge, say $ab$, then $(a,b)$ and $(b,a)$ are good pairs,
with triangles $adc$ and $bcd$. Unless $|cd|<|ad|,|ac|,|bc|,|bd|$, we are done.
In that remaining case $|ab|\le|cd|$, so $|ab|$ and $|cd|$ are both strictly
smaller than $|ac|,|ad|,|bc|,|bd|$, and $(a,b)$ and $(c,d)$ are good pairs.

Suppose that both searches get stuck after their first step. For the root
$a$, let $x\in(a,b]$ be the first point added; the witness is
$\{a,x,c,d\}$, and being stuck means: the next ccw point $X\in(x,c)$, the
first point after $x$ at distance at least $r$ from $a$ and $x$, satisfies
$|Xd|<r$; the next cw point $Y\in(d,a)$ satisfies $|Yc|<r$, and $|Ya|\ge r$.
Likewise for the root $c$, with first point $x'\in(c,d]$: the next ccw point
$X'\in(x',a)$ has $|X'c|\ge r$ and $|X'b|<r$, and the next cw point
$Y'\in(b,c)$ has $|Y'a|<r$.

First, $X\in(b,c)$: $X\ne b$ because $|bd|\ge r>|Xd|$, and if $X$ lay before
$b$, then (Q) for $a,X,b,d$ would give $|ab|>|bd|+(|aX|-|Xd|)>|bd|$, contrary
to $|ab|<|bd|$. Symmetrically $X'\in(d,a)$, using $|cd|<|bd|$. Second, $Y'$
precedes $X$ in $(b,c)$: otherwise $a,X,Y',c,d$ would contradict
\cref{lem:core}. Symmetrically, $Y$ precedes $X'$ in $(d,a)$ (apply
\cref{lem:core} to $c,X',Y,a,b$). Hence the eight points
$a,b,Y',X,c,d,Y,X'$ are in convex position in this order.

Apply \cref{lem:opposite} to $X',a,b,Y'$ (with $|ab|\ge r>|X'b|,|aY'|$) and to
$X,c,d,Y$ (with $|cd|\ge r>|Xd|,|cY|$): in these two quadrilaterals the angles
at $X'$ and $Y'$, and at $X$ and $Y$, sum to at least $\pi$ each. The octagon
is the union of these two quadrilaterals and the quadrilateral
$Y',X,Y,X'$, glued along $X'Y'$ and $XY$. Its interior angles at
$Y',X,Y,X'$ are therefore the sums of the corresponding angles of the pieces,
and they add up to at least $\pi+\pi+2\pi=4\pi$. But each is smaller than
$\pi$, a contradiction. So one of the two searches has a safe second step.
\end{proof}

Under the arc property the claim holds in a
strong form.

\begin{proposition}\label{prop:sbmarc}
If $P$ has the $r$-arc property (\cref{sec:arc}) and $T$ is an $r$-far
$k$-set, then from every root $q\in T$ the path of $k-1$ ccw steps
succeeds. In particular $(\mathrm{SBM}_k)$ holds, and no branching is
necessary.
\end{proposition}
\begin{proof}
By \cref{thm:arc}(i) the greedy chain $x_0=q$, $x_{i+1}=a(x_i)$ is proper and
$r$-far. We show by induction that the $i$-th ccw step adds $x_i$. The points
strictly between $x_{i-1}$ and $x_i$ are closer than $r$ to $x_{i-1}$, so they
are not eligible; $x_i$ is at distance at least $r$ from $x_0,\dots,x_{i-1}$
because the chain is $r$-far; and $x_i$ does not pass the cw end $q$ because
the chain is proper.
\end{proof}

\section{Natural choices of the root fail}
The proof of \cref{thm:sbm4} fixes the root in terms of $T$. The natural
candidates for such a rule do not work for larger $k$.

\begin{proposition}\label{prop:rootrules}
There are point sets in convex position, a radius $r$ and an $r$-far $k$-set
$T$ such that
\begin{enumerate}[label=(\alph*),itemsep=1pt,topsep=2pt]
\item ($k=5$ and $k=6$) the two-sided search fails from an endpoint of a
closest pair of $T$;
\item ($k=7$) no sequence of safe steps from the vertex of $T$ with the
smallest interior angle reaches $k$ points;
\item ($k=7$) the same holds for the vertex of smallest angle over all $r$-far
$k$-sets of $P$.
\end{enumerate}
In all these instances $(\mathrm{SBM}_k)$ holds.
\end{proposition}
\begin{proof}
The instances have between $10$ and $14$ points with integer coordinates and
are listed in \code{fpt/data/root_rule_instances.json}; each claim is verified
by exhaustive search in exact arithmetic (\code{fpt/check_roots.py}).
\end{proof}


\chapter{When greedy is exact: the arc property}\label{sec:arc}

The counterexamples of \cref{sec:greedy} are thin polygons in which the
distance from a point is far from unimodal along the hull. This section shows
that this is the only obstruction: if the set of points far from $q$ is always
contiguous, the simplest greedy rule is exact, and for unimodal point sets it
gives a near-linear algorithm.

For $r>0$ and $q\in P$ let $\mathrm{Far}_r(q)=\{x\in P\setminus\{q\}:|qx|\ge r\}$.
We say that $P$ has the \emph{$r$-arc property} if for every $q$ the set
$\mathrm{Far}_r(q)$ consists of consecutive points in the order
$p_{q+1},\dots,p_{q-1}$; then $\mathrm{Far}_r(q)$ is exactly the ccw arc from
$a(q)$ to $b(q)$. We call $P$ \emph{unimodal} if for every $q$ the sequence
$|p_qp_{q+t}|$, $t=1,\dots,n-1$, is strictly increasing and then strictly
decreasing, where the maximum may be attained at two consecutive values of
$t$. A unimodal set has the $r$-arc property for every $r$. Points on a circle
are unimodal, since the chord length is $2R\sin(\theta/2)$ for the ccw angle
$\theta\in(0,2\pi)$; the tie at the top occurs, for instance, in every regular
polygon with an odd number of vertices, which is why it has to be admitted.

\begin{lemma}[Exchange]\label{lem:a1}
Assume the $r$-arc property. Let $S$ be an $r$-far set and let $s,s'$ be
consecutive points of $S$ in ccw order. Then $(S\setminus\{s'\})\cup\{a(s)\}$
is $r$-far.
\end{lemma}
\begin{proof}
By \cref{lem:arcbasic}, $a(s)$ lies on the ccw arc from $s$ to $s'$, and
$|s\,a(s)|\ge r$. Let $t\in S\setminus\{s,s'\}$. The set $\mathrm{Far}_r(t)$
is a set of consecutive points in the order that starts and ends next to $t$,
and it contains $s$ and $s'$. The ccw arc from $s$ to $s'$ does not contain
$t$, so all its points lie between $s$ and $s'$ in that order and belong to
$\mathrm{Far}_r(t)$. Hence $|t\,a(s)|\ge r$.
\end{proof}

For $q\in P$ and $k\ge2$ the \emph{greedy chain} of $q$ is $x_0=q$,
$x_{i+1}=a(x_i)$ for $i<k-1$. It is \emph{proper} if all $a(x_i)$ are defined
and the chain does not wrap around, i.e.\ the ccw offsets of $x_0,\dots,x_{k-1}$
from $q$ are strictly increasing and smaller than $n$.

\begin{theorem}\label{thm:arc}
Assume the $r$-arc property and let $k\ge2$.
\begin{enumerate}[label=(\roman*),itemsep=1pt,topsep=2pt]
\item If $S$ is an $r$-far $k$-set and $q\in S$, then the greedy chain of $q$
is proper and $r$-far.
\item A proper greedy chain is $r$-far if and only if $|x_0x_{k-1}|\ge r$.
\end{enumerate}
Hence $P$ has an $r$-far $k$-set if and only if some $q$ has a proper greedy
chain with $|x_0x_{k-1}|\ge r$.
\end{theorem}
\begin{proof}
(i) Let $S=\{s_1,\dots,s_k\}$ in ccw order with $s_1=q$. Applying
\cref{lem:a1} to the consecutive pair $(x_0,s_2)$ replaces $s_2$ by $x_1$;
then $(x_1,s_3)$ is a consecutive pair of the new $r$-far set, and we replace
$s_3$ by $x_2$, and so on. After $k-1$ steps the set is the chain, which is
$r$-far, and since $x_i$ lies on the ccw arc from $x_{i-1}$ to $s_{i+1}$, the
chain is proper.

(ii) Only ``if'' needs a proof. Fix $i$. Starting after $x_i$, the other
chain points appear in the order
$x_{i+1},\dots,x_{k-1},x_0,\dots,x_{i-1}$. The first of them is $a(x_i)$ and
lies in $\mathrm{Far}_r(x_i)$. The last one also does: for $i\ge1$ because
$x_i=a(x_{i-1})$, and for $i=0$ by assumption. Since $\mathrm{Far}_r(x_i)$ is
contiguous, it contains all of them.
\end{proof}

Part~(ii) is the difference to \cref{sec:greedy}: without the arc property a
greedy chain can close up and still contain a close pair, and by
\cref{thm:tightpair} there may be no $r$-far chain at all. Indeed the set
$P_\star$ violates the arc property at $r^2=10240$ ($\mathrm{Far}_r(p_0)$
consists of the three runs $\{p_1,p_2\}$, $\{p_4\}$, $\{p_6\}$).

\paragraph{Decision.}
Given the offsets of $a(q)$ for all $q$, the function $a^{k-1}$ together with
the offset travelled is obtained by repeated squaring in $O(n\log k)$ time and
$O(n)$ space, after which \cref{thm:arc} is tested in $O(1)$ per seed. (The
obvious $O(nk)$ walk works as well.) Computing $a(q)$ takes $O(n)$ time by
scanning. For unimodal $P$ we first find, for every row $q$, the position of
its maximum by binary search; then $a(q)$ is found by binary search on the
increasing part of the row.

\begin{theorem}\label{thm:arcdecision}
Under the $r$-arc property, the decision problem can be solved in
$O(n^2)$ time, and in $O(n\log n)$ time if $P$ is unimodal. A witness is
reported within the same bounds.
\end{theorem}

\paragraph{Optimisation.}
For unimodal $P$, sorting the $\binom n2$ distances and binary search gives
$O(n^2\log n)$ time; we call this variant \textsc{O1}. To avoid generating
all distances, observe that every row splits at its maximum into an increasing
and a decreasing list, so that the candidates are available as $2n$ implicit
sorted lists. Algorithm \textsc{O2} maintains an open interval $(lo,hi)$ with
$lo$ feasible and $hi$ infeasible and, for each list, the range of its entries
inside the interval. In each round it picks one of the remaining entries
uniformly at random, tests it, and shrinks all ranges by binary search.

\begin{theorem}\label{thm:arcopt}
For $n$ points in convex position that form a unimodal set, $k$-dispersion can
be solved exactly in expected $O(n\log^2n)$ time and $O(n)$ space.
\end{theorem}
\begin{proof}
A round costs $O(n\log n)$: $O(n)$ to select the pivot, $O(n\log n)$ for the
decision (\cref{thm:arcdecision}) and $O(n\log n)$ to update the ranges. All
remaining entries equal to the pivot are removed, and with probability at
least $1/2$ the pivot has rank between $1/4$ and $3/4$ among the remaining
entries, in which case at least a quarter of them are removed. So the expected
number of rounds is $O(\log n)$. When no entry is left, $lo=\rstar$, because
$\rstar$ is a pairwise distance and is feasible.
\end{proof}

\paragraph{Recognising the hypotheses.}
The arc property at a given $r$ is checked in $O(n^2)$ time by counting runs.
Unimodality can be verified faster.

\begin{lemma}\label{lem:unitest}
Whether $P$ is unimodal can be decided in $O(n\log n)$ time.
\end{lemma}
\begin{proof}
Fix $j$. The points strictly closer to $p_j$ than to $p_{j+1}$ are the points
of $P$ in an open half-plane bounded by the perpendicular bisector of
$p_jp_{j+1}$. Since $P$ is in strictly convex position, they form a run of
consecutive points; it contains $p_j$ and not $p_{j+1}$, so it is
$p_j,p_{j-1},\dots,p_{j-c_j+1}$ for some $1\le c_j\le n-1$. At most two points
of $P$ lie on the bisector; one boundary crossing of the bisector is inside the
edge $p_jp_{j+1}$, so the only candidate for a point equidistant from $p_j$ and
$p_{j+1}$ is $p_{j-c_j}$. Let $z_j=1$ if it is equidistant and $z_j=0$
otherwise. The predicate ``$p_{j-s}$ is strictly closer to $p_j$ than to
$p_{j+1}$'' is monotone in $s$, so $c_j$ is found by binary search with
$O(\log n)$ distance comparisons, and $z_j$ with one more.

Now fix a row $q$ and $1\le t\le n-2$, and put $j=q+t$. The step from
$p_{q+t}$ to $p_{q+t+1}$ increases the distance to $p_q$ iff $p_q=p_{j-t}$ is
strictly closer to $p_j$ than to $p_{j+1}$, that is, iff $t<c_j$; it leaves the
distance unchanged iff $t=c_j$ and $z_j=1$; otherwise it decreases. A row is
unimodal (with at most one tie, at the top) iff no step that is not an
increase is followed by a step that is not a decrease. Step $t$ is not an
increase iff $t\ge c_j$, and step $t+1$ (whose edge is $p_{j}p_{j+1}$ shifted
by one, i.e.\ index $j+1$) is not a decrease iff $t+1\le c_{j+1}-1+z_{j+1}$.
Hence some row fails iff for some $j$ there is an integer
$t\in[1,n-3]$ with $c_j\le t\le c_{j+1}+z_{j+1}-2$. As $c_j\ge1$, this is the
case iff $c_j\le\min\{c_{j+1}+z_{j+1}-2,\;n-3\}$ for some $j$, which is
checked in $O(n)$ time once all $c_j,z_j$ are known.
\end{proof}

For general convex input we use an \emph{adaptive} variant of \textsc{O1}: it
checks the arc property at every radius it tests and reports ``not
applicable'' if the check fails. It takes $O(n^2\log n)$ time, is exact
whenever it answers, and never answers incorrectly.

\paragraph{Weaker hypotheses.}
The hypothesis of \cref{thm:arc} cannot be weakened in the two obvious ways.

\begin{proposition}\label{prop:weaker}
\begin{enumerate}[label=(\alph*),itemsep=1pt,topsep=2pt]
\item There is a set of $10$ points and an $r$-far $6$-set $S$ such that
$\mathrm{Far}_r(s)$ is contiguous for every $s\in S$, but no greedy chain of six
points is $r$-far.
\item There is a set of $6$ points with the $\rstar$-arc property for $k=4$ on
which binary search over the candidate distances with the test of
\cref{thm:arc}(ii) returns a value larger than $\rstar$.
\end{enumerate}
\end{proposition}
\begin{proof}
The instances are listed in \code{arc/data/hypothesis_counterexamples.json} and
verified by exhaustive search in exact arithmetic
(\code{arc/check\_hypotheses.py}).
\end{proof}

Running times of the near-linear algorithm are reported in \cref{app:arcexp}.


\chapter{A Simple Exact Algorithm}\label{ch:exact}
\section{The capped Delaunay dynamic program}\label{sec:dp}

Fix $r>0$ and write $\far$ for $\far_r$. For the decision problem it suffices
to consider \emph{chords} $(i,j)$ with $0\le i<j\le n-1$. A chord is
\emph{admissible} if $\far(i,j)$, and a point $p_l\in P(i,j)$ is an
\emph{admissible apex} of $(i,j)$ if $\far(i,l)$ and $\far(l,j)$. Let
\[
\Tr=\bigl|\{(i,l,j): 0\le i<l<j\le n-1,\ \far(i,l),\ \far(l,j),\ \far(i,j)\}\bigr|\le\tbinom n3
\]
be the number of $r$-far triples.

\begin{definition}\label{def:config}
Let $(i,j)$ be an admissible chord. An \emph{$(i,j)$-configuration} is a pair
$(X,T)$ where $X\subseteq P(i,j)$ and, if $X\neq\emptyset$, $T$ is a
triangulation of $X\cup\{p_i,p_j\}$ that satisfies (a) and (b) of
\cref{thm:cert-delaunay}. If $X\ne\emptyset$, the \emph{apex} of $(X,T)$ is
the third vertex $p_l$ of the triangle of $T$ containing the edge $p_ip_j$.
For $\varphi\in(0,\pi]$ let $F(i,j\mid\varphi)$ be the maximum of $|X|$ over
all $(i,j)$-configurations with $X=\emptyset$ or $\ang p_ip_lp_j\le\varphi$.
\end{definition}

The bound $\varphi$ encodes the triangle on the other side of the chord: if
$p_ip_j$ is a diagonal of a larger triangulation whose other triangle at
$p_ip_j$ is $p_ip_jp_t$, then by~\eqref{eq:ld} the diagonal is locally
Delaunay exactly if $\ang p_ip_lp_j\le\pi-\ang p_ip_tp_j$. See
\cref{fig:dp}.

\begin{figure}[tb]
\centering
\begin{tikzpicture}[scale=1.0,>=stealth]
  \coordinate (i) at (-2.6,0);
  \coordinate (j) at (2.6,0);
  \coordinate (l) at (0.5,-2.2);
  \coordinate (x) at (-1.75,-1.75);
  \coordinate (y) at (2.05,-1.45);
  \coordinate (t) at (0.2,1.9);
  % circle through i, l, j: centre (0, c) with c = (0.25+4.84-6.76)/(-4.4)
  \draw[gray!70,dashed] (0,0.3795) circle (2.6276);
  \fill[blue!8] (i)--(l)--(j)--cycle;
  \draw[thick] (i)--(j);
  \draw[thick] (i)--(l)--(j);
  \draw (i)--(x)--(l);
  \draw (l)--(y)--(j);
  \draw[gray] (i)--(t)--(j);
  \foreach \p/\pos/\lab in {i/left/{$p_i$}, j/right/{$p_j$}, l/below/{$p_l$}, x/below left/{$x=L_c(i,l)$}, y/below right/{$y=L_{c'}(l,j)$}, t/above/{$p_t$}}
    {\fill (\p) circle (1.6pt); \node[\pos] at (\p) {\lab};}
  \node at (0.15,-0.75) {\small apex triangle};
  \node[gray] at (2.75,2.25) {\small $\circlethree{p_i}{p_l}{p_j}$};
\end{tikzpicture}
\caption{A chord $(i,j)$ with apex $p_l$. The sub-configuration of $(i,l)$
with apex $x$ is compatible with the triangle $p_ip_lp_j$ iff $x$ is not
strictly inside $\circlethree{p_i}{p_l}{p_j}$ (here: $x$ is outside); likewise
for $y$ and $(l,j)$. The parent triangle $p_ip_jp_t$ constrains $p_l$ in the
same way.}
\label{fig:dp}
\end{figure}

\begin{lemma}[Recurrence]\label{lem:rec}
For every admissible chord $(i,j)$ and every $\varphi\in(0,\pi]$,
\[
F(i,j\mid\varphi)=\max\Bigl(\{0\}\cup\bigl\{\,v(l) : p_l\text{ admissible apex of }(i,j),\ \ang p_ip_lp_j\le\varphi\bigr\}\Bigr),
\]
where
$v(l)=1+F\bigl(i,l\mid\pi-\ang p_ip_jp_l\bigr)+F\bigl(l,j\mid\pi-\ang p_lp_ip_j\bigr)$.
\end{lemma}
\begin{proof}
($\ge$) Let $p_l$ be an admissible apex with $\ang p_ip_lp_j\le\varphi$, and
let $(X_1,T_1)$ and $(X_2,T_2)$ be optimal configurations of $(i,l)$ and
$(l,j)$ under the bounds in the definition of $v(l)$. Put
$X=X_1\cup X_2\cup\{p_l\}$ and let $T$ consist of $T_1$, $T_2$ and the
triangle $p_ip_lp_j$. The three sides of this triangle have length at least
$r$. If $X_1\ne\emptyset$, then $p_ip_l$ is a diagonal of $T$ whose two
triangles are $p_ip_lp_j$ and the apex triangle of $T_1$, and
by~\eqref{eq:ld} the bound $\pi-\ang p_ip_jp_l$ on the apex angle of $T_1$ is
exactly its local Delaunay condition; symmetrically for $p_lp_j$. All other
diagonals of $T$ are diagonals of $T_1$ or $T_2$. So $(X,T)$ is an
$(i,j)$-configuration with apex $p_l$ and $|X|=v(l)$.

($\le$) Let $(X,T)$ be an optimal configuration with $X\ne\emptyset$ and apex
$p_l$. By (a), $p_l$ is an admissible apex. Removing the triangle $p_ip_lp_j$
splits $T$ into triangulations of $(X\cap P(i,l))\cup\{p_i,p_l\}$ and
$(X\cap P(l,j))\cup\{p_l,p_j\}$, which are configurations of $(i,l)$ and
$(l,j)$, and by (b) their apex angles obey the bounds in $v(l)$. Hence
$|X|\le v(l)$.
\end{proof}

\begin{theorem}[Decision]\label{thm:decision}
Let $k\ge3$. $P$ has an $r$-far $k$-set if and only if some admissible chord
$(i,j)$ has $F(i,j\mid\pi)\ge k-2$.
\end{theorem}
\begin{proof}
($\Rightarrow$) Let $S$ be an $r$-far $k$-set, let $p_i$ and $p_j$ be its
points of smallest and largest index, and let $T$ be a Delaunay triangulation
of $S$. Then $X=S\setminus\{p_i,p_j\}\subseteq P(i,j)$ and $(X,T)$ is an
$(i,j)$-configuration by \cref{thm:cert-delaunay}, with $|X|=k-2$.
($\Leftarrow$) A configuration $(X,T)$ with $|X|\ge k-2\ge 1$ is a
triangulation of $X\cup\{p_i,p_j\}$ satisfying (a) and (b). By
\cref{thm:cert-delaunay} this set is $r$-far, and any $k$ of its points form
an $r$-far $k$-set.
\end{proof}

\subsection{Thresholds: \texorpdfstring{$k-2$}{k-2} apex points per chord}

Evaluating the recurrence naively needs one table entry per chord and apex and
$O(n)$ work per entry, that is, $O(n^4)$ time. The following observation
removes a factor~$n$ and, with capping, replaces $\log n$ by $\log k$.

Fix an admissible chord $(i,j)$. For $c\ge1$ define the \emph{$c$-th
threshold}
\[
\tau_c(i,j)=\min\bigl\{\ang p_ip_lp_j : p_l\text{ admissible apex of }(i,j),\ v(l)\ge c\bigr\},
\]
with $\tau_c(i,j)=+\infty$ if no such apex exists, and put $\tau_0=0$.

\begin{lemma}\label{lem:thresholds}
For every admissible chord, $\tau_1\le\tau_2\le\cdots$ and
$F(i,j\mid\varphi)=\max\{c\ge0:\tau_c(i,j)\le\varphi\}$. Moreover, for every
integer $K\ge1$, the thresholds $\tau_1,\dots,\tau_K$ of a chord are determined
by the thresholds $\tau_1,\dots,\tau_K$ of shorter chords.
\end{lemma}
\begin{proof}
The sets $\{l:v(l)\ge c\}$ shrink as $c$ grows, which gives monotonicity. By
\cref{lem:rec}, $F(i,j\mid\varphi)\ge c$ iff some admissible apex $p_l$ with
$\ang p_ip_lp_j\le\varphi$ has $v(l)\ge c$, iff $\tau_c\le\varphi$. For the
last claim write $F_K=\min(F,K)$. The thresholds $\tau_1,\dots,\tau_K$ depend
on $v$ only through $\min(K,v)$, and
$\min(K,v(l))=\min\bigl(K,1+F_K(i,l\mid\cdot)+F_K(l,j\mid\cdot)\bigr)$, where
$F_K$ of a chord is determined by its first $K$ thresholds.
\end{proof}

For a dispersion query we only need to know whether $F\ge k-2$, so we set
$K=k-2$ and keep $K$ thresholds per chord. We do not store angles: we store
$\tau_c(i,j)$ as an apex $L_c(i,j)$ that attains it ($L_c=\nil$ if
$\tau_c=+\infty$). Then, for an admissible apex $p_l$ of $(i,j)$ and
$x=L_c(i,l)$,
\[
F_K\bigl(i,l\mid\pi-\ang p_ip_jp_l\bigr)\ge c
\iff \ang p_ixp_l+\ang p_ip_jp_l\le\pi
\iff x\text{ is not strictly inside }\circlethree{p_i}{p_l}{p_j}
\]
by (I2), since $x$ and $p_j$ lie on opposite sides of the line $p_ip_l$. The
same circle answers the query for the chord $(l,j)$. Each query is monotone in
$c$, so the number of thresholds that pass is found by binary search with
$O(\log K)$ in-circle tests. Comparing the apex angles of two apexes $p_l$,
$p_{l'}$ of the same chord is the in-circle test~(I1). \Cref{alg:capped}
gives the complete procedure.

\begin{algorithm}
\caption{$\AlgCapped(P,k,r)$\quad ($k\ge3$, $K=k-2$)}\label{alg:capped}
\begin{tabbing}
\hspace{1.4em}\=\hspace{1.4em}\=\hspace{1.4em}\=\hspace{1.4em}\=\kill
\textbf{for} $g=2,\dots,n-1$ \textbf{do}\` \emph{chords by increasing length $j-i$}\\
\>\textbf{for each} admissible chord $(i,j)$ with $j=i+g$ \textbf{do}\\
\>\>$best[1..K]\gets\nil$\\
\>\>\textbf{for each} admissible apex $p_l$ of $(i,j)$ \textbf{do}\` \emph{bit-parallel enumeration}\\
\>\>\>$q_1\gets\#\{c: L_c(i,l)\ne\nil\text{ is not strictly inside }\circlethree{p_i}{p_l}{p_j}\}$\` \emph{binary search}\\
\>\>\>$q_2\gets\#\{c: L_c(l,j)\ne\nil\text{ is not strictly inside }\circlethree{p_i}{p_l}{p_j}\}$\` \emph{binary search}\\
\>\>\>$v\gets\min(K,1+q_1+q_2)$\\
\>\>\>\textbf{if} $best[v]=\nil$ or $\ang p_ip_lp_j<\ang p_i\,best[v]\,p_j$ \textbf{then} $best[v]\gets p_l$\\
\>\>\textbf{for} $c=K,\dots,1$ \textbf{do} $L_c(i,j)\gets$ apex of smallest angle among $best[c..K]$\` \emph{suffix minima}\\
\>\>\textbf{if} $L_K(i,j)\ne\nil$ \textbf{then return} \textsc{true}\\
\textbf{return} \textsc{false}
\end{tabbing}
\end{algorithm}

\begin{theorem}\label{thm:decision-alg}
\Cref{alg:capped} decides whether $P$ has an $r$-far $k$-set in
\[
O\bigl(n^2k+n^3/w+\Tr\log k\bigr)\ \subseteq\ O(n^3\log k)
\]
time and $O(n^2k)$ space. If the answer is positive, an $r$-far $k$-set can be
reported in $O(k\log k)$ additional time.
\end{theorem}
\begin{proof}
\emph{Correctness.} By induction on $j-i$, after the chord $(i,j)$ has been
processed, $L_c(i,j)$ attains $\tau_c(i,j)$ for $1\le c\le K$: the counts
$q_1,q_2$ equal $F_K$ of the two sub-chords under the bounds of
\cref{lem:rec} by the discussion above, so $v=\min(K,v(l))$; $best[v]$ is an
apex of smallest angle among the apexes with capped value exactly $v$, and the
suffix minima give the apexes of smallest angle with capped value at least
$c$. By \cref{lem:thresholds} and \cref{thm:decision}, an $r$-far $k$-set
exists iff $\tau_K(i,j)\le\pi$, i.e.\ $L_K(i,j)\ne\nil$, for some chord.

\emph{Running time.} Initialising and finishing a chord costs $O(K)$. The
admissible apexes of $(i,j)$ are the positions $l\in(i,j)$ where the rows of
$p_i$ and $p_j$ in the bit matrix of the relation $\far$ both have a one; they
are enumerated in $O(n/w)$ time plus $O(1)$ per apex. The body of the inner
loop is executed once per $r$-far triple and costs $O(\log K)$.

\emph{Witness.} Starting from a chord with $L_K\ne\nil$ and demand $c=K$,
take $p_l=L_c(i,j)$, recompute $q_1$ for $p_l$, and recurse on $(i,l)$ with
demand $c_1=\min(q_1,c-1)$ and on $(l,j)$ with demand $c-1-c_1\le q_2$. The
recursion visits at most $K$ chords and performs one binary search in each.
The reported points together with $p_i,p_j$ form the vertex set of a
configuration and are $r$-far by \cref{thm:cert-delaunay}.
\end{proof}

\begin{remark}[The role of the cap]\label{rem:cap}
With $K=n-2$ nothing is truncated and $2+\max_{(i,j)}F(i,j\mid\pi)$ is the
maximum size of an $r$-far subset, in $O(n^3\log n)$ time and $O(n^3)$ space;
this matches the bound of~\citet{tw26} in the unweighted unit case. Capping
reduces the logarithm to $\log k$, the space to $O(n^2k)$, and allows the
early exit in the last line of the loop. The $\log k$ factor is only paid on
$r$-far triples; in the upper part of the search window many triples are not
$r$-far.
\end{remark}

\paragraph{Maximum independent sets.}
The cap can also be chosen adaptively. Let $\alpha=\alpha_r(P)$ be the maximum
size of an $r$-far subset of $P$, i.e.\ the independence number of the graph on
$P$ in which two points are adjacent iff their distance is smaller than $r$
(the unit-disk graph, after scaling).

\begin{corollary}\label{cor:mis}
A maximum $r$-far subset of $n$ points in convex position can be computed in
$O(n^2\alpha+n^3/w+\Tr\log\alpha)\subseteq O(n^3\log\alpha)$ time and
$O(n^2\alpha)$ space, where $\alpha$ is its size.
\end{corollary}
\begin{proof}
Compute an inclusion-maximal $r$-far set $M$ greedily in $O(n|M|)$ time. Every
point of $P$ is at distance less than $r$ from some point of $M$, and an open
disk of radius $r$ contains at most five points of any $r$-far set (among six
points other than the centre, two would subtend an angle of at most $\pi/3$ at
the centre and would therefore be closer than $r$; and if the centre belongs
to the set, no other point of the disk does).
Hence $|M|\le\alpha\le5|M|$. Run \cref{alg:capped} with
$K=\min(n,5|M|)-1$ and without the early exit. Since $F\le\alpha-2<K$, no
value is truncated, so the maximum number of finite thresholds over all chords
equals $\alpha-2$ if $\alpha\ge3$, and a witness is reconstructed as in
\cref{thm:decision-alg}. If there is no $r$-far triple, the answer is an
$r$-far pair or a single point.
\end{proof}

\section{Optimisation and the search window}\label{sec:opt}

The optimum $\rstar$ is one of the $\binom n2$ pairwise distances, and
feasibility is monotone in $r$, so $\rstar$ can be found by binary search over
the sorted distinct distances with \cref{alg:capped} as the oracle. The
following window shrinks the search range before any exact decision is made.

\begin{lemma}[Search window]\label{lem:window}
Let $G$ be a farthest-first $k$-set of $P$ (start with an arbitrary point and
repeatedly add a point whose distance to the chosen points is largest), and
let $g=\delta(G)$. Then for $k\ge3$
\[
g\ \le\ \rstar\ \le\ \min\bigl\{2g,\ \per(P)/k\bigr\}.
\]
\end{lemma}
\begin{proof}
$g\le\rstar$ because $G$ is a feasible $k$-set. Let $G'$ be the set of the
first $k-1$ points chosen. The distances at which points are added are
non-increasing, so $g$ is the distance from the last point to $G'$, and every
point of $P$ is within distance $g$ of $G'$ (this is Gonzalez's
argument~\citep{gonzalez}). Two of the $k$ points of an optimal set have the
same nearest point in $G'$, hence $\rstar\le2g$. Finally, let $S^*$ be an
optimal set. Then $\conv(S^*)$ is a convex $k$-gon whose sides have length at
least $\rstar$, so its perimeter is at least $k\rstar$; since
$\conv(S^*)\subseteq\conv(P)$ and the perimeter of convex sets is monotone
under inclusion, $k\rstar\le\per(P)$.
\end{proof}

\begin{theorem}\label{thm:opt}
The $k$-dispersion problem for $n$ points in convex position can be solved in
$O(n^3\log k\log n)$ time and $O(n^2k)$ space. More precisely, the time is
$O(n^2\log n)$ plus $O(\log n)$ calls to \cref{alg:capped}.
\end{theorem}
\begin{proof}
For $k=2$ the answer is the diameter. For $k\ge3$, sort the distinct pairwise
distances, restrict them to the window of \cref{lem:window} (computing $G$
takes $O(nk)$ time), and binary search with \cref{alg:capped}; the value $g$
is known to be feasible. The witness of the last positive call, or $G$ if there
is none, is an optimal set.
\end{proof}

\subsection{A certified algorithm}\label{sec:certified}

Although no seeded rule is complete, all of them are \emph{sound}: they only
ever output $r$-far sets. So a successful greedy run at radius $\rho$ proves
$\rstar\ge\rho$, and a single exact decision just above $\rho$ can prove
$\rstar\le\rho$. \Cref{alg:certified} combines the two. As heuristic predicate
$H(r)$ we use the alternating greedy from all seeds $\{q,a(q)\}$ and
$\{q,a(q),b(q)\}$ with both starting sides.

\begin{algorithm}[H]
\caption{$\AlgCert(P,k)$\quad ($k\ge3$)}\label{alg:certified}
\begin{enumerate}[label=\arabic*:,leftmargin=2em,itemsep=1pt,topsep=2pt]
\item $W\gets$ the sorted distinct distances inside the window of \cref{lem:window};\ \ $\rho\gets g$, which is feasible.
\item Binary search on $W$ with the predicate $H$; let $\rho$ be the largest radius at which a run succeeded.
\item Let $\rho^+$ be the successor of $\rho$ in $W$. \textbf{If} $\rho=\max W$ \textbf{or} $\AlgCapped(P,k,\rho^+)$ is \textsc{false}, \textbf{return} $\rho$.
\item Otherwise $\rho^+$ is feasible: \textbf{return} the result of binary search on $W\cap[\rho^+,\infty)$ with $\AlgCapped$.
\end{enumerate}
\end{algorithm}

\begin{theorem}\label{thm:cert}
\Cref{alg:certified} returns $\rstar$ on every input. It runs in
$O(n^2k\log n+n^3\log k)$ time whenever $\rho=\rstar$ after the heuristic
phase, and in $O(n^3\log k\log n)$ time in the worst case.
\end{theorem}
\begin{proof}
Every success of $H$ exhibits an $r$-far $k$-set and $g$ is feasible, so
$\rho\le\rstar$ regardless of whether $H$ is monotone. If $\rho=\max W$, then
$\rho=\rstar$ by \cref{lem:window}. If the exact decision at $\rho^+$ fails,
then $\rstar<\rho^+$ and hence $\rstar=\rho$, because $\rstar\in W$. Otherwise
$\rstar\ge\rho^+$ and the exact binary search finds it.

For the running time, $a(q)$ and $b(q)$ are found for all $q$ in $O(n^2)$
time. In one greedy run each frontier only moves forward, and a point that is
skipped is never examined again from that side, so a run examines $O(n)$
points and compares each with at most $k$ chosen points. There are $4n$ runs
per evaluation of $H$ and $O(\log n)$ evaluations, followed by a single call
to \cref{alg:capped} if $\rho=\rstar$.
\end{proof}


\chapter{Implementation and Results}\label{ch:results}

\begin{chapterabstract}
This section describes how the algorithms were implemented and validated,
and reports the experimental results: correctness checks against exhaustive
search, running times of the exact algorithms, the effect of the cap and of the
search window, the certified algorithm, and the near-linear algorithm for
unimodal point sets.
\end{chapterabstract}

\section*{Methodology}
\addcontentsline{toc}{section}{Methodology}

The project followed a cycle of \emph{conjecture, test, prove}. Every structural
property studied in \cref{sec:greedy,sec:sbm,sec:arc} was first implemented as a
predicate and tested on large numbers of random and adversarially optimised
instances against an exhaustive solver. Properties that failed produced
explicit counterexamples, which were then minimised and verified independently
in Python with exact integer arithmetic; properties that survived were proved.

The code is organised in four directories of the repository:
\code{convex/} (the general algorithms and their benchmarks), \code{arc/}
(the arc-property algorithms), \code{sbm/} (a re-implementation of the
algorithms of~\citet{sbm}) and \code{fpt/} (the tools used to analyse the
search tree).

\section{Experiments with the general algorithms}\label{sec:experiments}

\paragraph{Implementation.}
All algorithms are implemented in C++17 in a single header
(\code{convex/kdisp.hpp}); geometric decisions use only comparisons of $64$-bit
squared distances and the $128$-bit in-circle determinant. Besides
\cref{alg:capped,alg:certified} the code contains two baselines: the same
dynamic program without the cap ($K=n-2$, cf.\ \cref{rem:cap}), and the direct
\emph{triple-state} program that stores a value for every chord and apex and
scans all sub-apexes, which takes $O(n^4)$ in-circle tests per decision. All
optimisation variants use the search window unless stated otherwise.
Experiments were run on one core of a virtual machine (x86-64, GCC~13.3,
\code{-O2}); times are medians of three runs for $n\le200$ and single runs
for larger instances, and exclude reading the input and building the distance
matrix.

\paragraph{Instances.}
Every instance is generated by sampling a cloud of random integer points,
taking its strictly convex hull, and keeping $n$ random hull vertices. The
distributions are: uniform in a square and Gaussian (small hulls only),
uniform in a disk, uniform in a thin annulus, uniform in a $1{:}20$ ellipse
(\code{thin}), rounded points on a circle, exactly cocircular lattice points
(the $972$ lattice points on a circle of radius $1185665$, a worst case for
in-circle tests), and a thin lens bounded by two flat parabolic arcs
(\code{twochain}), the shape on which greedy rules fail. Coordinates are
bounded by $1.2\cdot10^{6}$. There are four splits: \emph{verify}
($6\le n\le14$, $\nVerify$ instances), \emph{medium} ($20\le n\le40$,
$\nMedium$ instances), \emph{scaling} ($25\le n\le400$, $k\in\{4,8,16,32\}$,
$\nScaling$ instances), and a sweep over $k$ at $n=200$.

\paragraph{Correctness.}
On every instance of every split, the capped program with and without the
window, the uncapped program, the certified algorithm and (for $n\le100$) the
triple-state program returned the same optimum, and every reported set was
checked to be a $k$-set with exactly that minimum distance
(\cref{tab:correct}). On \emph{verify} the optimum also agrees with a
branch-and-bound over all $k$-subsets, with an independent brute force
written in Python ($\nVerifyPy/\nVerify$), and with an independent
general-dimension solver based on clique enumeration and matrix products that
is part of the same code base ($\nVerifyFmm/\nVerify$). On \emph{medium},
exhaustive search was run for $k\le5$ ($\nMediumBrute$ instances). The
explicit instances of \cref{sec:greedy} are re-verified by a separate script
that shares no code with the solvers. In addition, a randomised self-test
compares all solvers and \cref{cor:mis} with exhaustive search on tens of
thousands of small instances, including degenerate ones with many equal
distances and cocircular quadruples.

\begin{table}[tb]
\centering
\caption{Cross-validation. ``agree'': all solvers return the same optimum and
valid optimal sets; ``exhaustive'': number of instances on which this value
was also confirmed by exhaustive search.}
\label{tab:correct}
\small
\begin{tabular}{@{}llrrrr@{}}
\toprule
split & distribution & instances & $n$ & agree & exhaustive\\
\midrule
\tabinput{generated/table_correctness}
\bottomrule
\end{tabular}
\end{table}

\paragraph{Running time.}
\Cref{fig:time} shows the median running times on the
\emph{scaling} split. The triple-state baseline grows roughly like $n^4$ and is
already $\tripleSpeedupHundred$ times slower than the capped program at
$n=100$.
At the optimum, a median of $\farTripleMedian\%$ of all triples are
$r$-far (between $\farTripleMin\%$ and $\farTripleMax\%$), so the term
$\Tr\log k$ dominates in practice.

\begin{figure}[tb]
\centering
\begin{tikzpicture}
\begin{groupplot}[group style={group size=4 by 1, horizontal sep=0.75cm},
  width=0.285\textwidth, height=0.36\textwidth, xmode=log, ymode=log,
  xlabel={$n$}, xtick={25,50,100,200,400}, xticklabels={25,50,100,200,400}, x tick label style={font=\tiny},
  tick label style={font=\scriptsize}, label style={font=\small},
  title style={font=\small}, grid=major, grid style={gray!25},
  legend style={font=\scriptsize, at={(2.55,1.22)}, anchor=south, legend columns=4, draw=gray!50, /tikz/every even column/.append style={column sep=8pt}},
  every axis plot/.append style={mark size=1.5pt}]
\nextgroupplot[title={$k=4$}, ylabel={time [ms]}]
  \addplot[black, mark=triangle*] table {generated/plot_k4_triple.dat};
  \addplot[orange!85!black, mark=diamond*] table {generated/plot_k4_uncapped.dat};
  \addplot[blue, mark=*] table {generated/plot_k4_capped.dat};
  \addplot[red!80!black, mark=square*] table {generated/plot_k4_certified.dat};
  \legend{triple-state, uncapped, capped, certified}
\nextgroupplot[title={$k=8$}]
  \addplot[black, mark=triangle*] table {generated/plot_k8_triple.dat};
  \addplot[orange!85!black, mark=diamond*] table {generated/plot_k8_uncapped.dat};
  \addplot[blue, mark=*] table {generated/plot_k8_capped.dat};
  \addplot[red!80!black, mark=square*] table {generated/plot_k8_certified.dat};
\nextgroupplot[title={$k=16$}]
  \addplot[black, mark=triangle*] table {generated/plot_k16_triple.dat};
  \addplot[orange!85!black, mark=diamond*] table {generated/plot_k16_uncapped.dat};
  \addplot[blue, mark=*] table {generated/plot_k16_capped.dat};
  \addplot[red!80!black, mark=square*] table {generated/plot_k16_certified.dat};
\nextgroupplot[title={$k=32$}]
  \addplot[black, mark=triangle*] table {generated/plot_k32_triple.dat};
  \addplot[orange!85!black, mark=diamond*] table {generated/plot_k32_uncapped.dat};
  \addplot[blue, mark=*] table {generated/plot_k32_capped.dat};
  \addplot[red!80!black, mark=square*] table {generated/plot_k32_certified.dat};
\end{groupplot}
\end{tikzpicture}
\caption{Median time to solve $k$-dispersion exactly, over all distributions
of the \emph{scaling} split (log--log). The triple-state baseline was run for
$n\le100$ and the uncapped program for $n\le200$.}
\label{fig:time}
\end{figure}

\paragraph{Effect of the cap.}
\Cref{tab:ksweep} compares the capped and the uncapped program at $n=200$.

\begin{table}[tb]
\centering
\caption{Sweep over $k$ at $n=200$ (median over six distributions): time [ms]
of the uncapped and the capped program and of the certified algorithm, and
in-circle tests per exact decision of the capped program (in millions).}
\label{tab:ksweep}
\small
\begin{tabular}{@{}rrrrrr@{}}
\toprule
$k$ & uncapped & capped & uncapped/capped & certified & in-circle tests [$10^6$]\\
\midrule
\tabinput{generated/table_ksweep}
\bottomrule
\end{tabular}
\end{table}

\paragraph{Certification.}
\Cref{tab:certified} reports how often the heuristic phase of
\cref{alg:certified} already attains the optimum, over all
$\certTotal$ instances of the verify, medium and scaling splits. This happens
on $\certYes$ instances, and on every instance drawn from the square,
Gaussian, disk, annulus, circle and cocircular distributions. In these cases
a single exact decision is made, and the certified algorithm is
$\certSpeedup$ times faster than the capped binary search in the median
(interquartile range $\certSpeedupLo$--$\certSpeedupHi$). The farthest-first
set alone is optimal far less often, so the seeded greedies do most of the
work. The output is exact in every case.

\begin{table}[tb]
\centering
\caption{Quality of the heuristic lower bounds. ``FF optimal'': the
farthest-first set attains $\rstar$; ``certified'': the heuristic phase of
\cref{alg:certified} attains $\rstar$, so that one exact decision suffices.
The last two columns give the mean and the maximum relative gap
$1-g/\rstar$ of the farthest-first value $g$.}
\label{tab:certified}
\small
\begin{tabular}{@{}lrrrrrr@{}}
\toprule
distribution & instances & FF optimal & certified & rate & mean gap & max gap\\
\midrule
\tabinput{generated/table_certified}
\bottomrule
\end{tabular}
\end{table}


\section{Experiments for the arc property}\label{app:arcexp}

\paragraph{Verification.}
The implementation (\code{arc/arcdisp.hpp}) was validated before any timing.
On $\arcSmallInst$ instances with $n\le16$ from ten families, at every
candidate radius $r$ (and at $r^2+1$) at which the arc property holds, three
implementations of the decision procedure agree with exhaustive search for
all $k=2,\dots,n$ ($\arcDchecks$ comparisons), and \textsc{O1}, \textsc{O2}
and the adaptive variant return the exhaustive optimum with a valid witness
($\arcOptChecks$ instance--$k$ pairs). On $\arcBigInst$ instances with
$20\le n\le150$ they agree with \cref{alg:capped} ($\arcBigChecks$ checks).
An independent Python script recomputes the optima of exported instances by
brute force and checks \cref{thm:arc} directly ($\arcPyAtwo$ checks). No
discrepancy occurred.

\paragraph{Running time.}
Unimodal instances are points at jittered equally spaced angles on a circle,
rounded to integers; every instance is verified to be strictly convex and
unimodal by the exact $O(n\log n)$ test. With coordinates bounded by $2^{29}$
such instances exist for $n\le50000$; for $n\ge10^5$ we use a circle of radius
$2^{60}$ and $128$-bit squared distances. \textsc{O1} needs $\Theta(n^2)$
memory and was run for $n\le10^4$. Every reported witness was
checked by an exact closest-pair computation. \Cref{tab:arcscaling} shows the
medians over three instances. \textsc{O2} solves instances with a million
points in about $\arcMillion$ seconds, and its running time grows like
$n^{\arcSlope}$ between $n=10^3$ and $n=10^6$, compared with
$n^{\arcSlopeOone}$ for \textsc{O1}. At $n=400$ it is between $300$ ($k=4$) and $40\,000$ ($k=64$) times faster
than the cubic dynamic program.

\begin{table}[tb]
\centering
\caption{Median running time [ms] on unimodal instances (three instances per
cell). \textsc{DP}: \cref{alg:capped} with binary search. ``--'': not run.}
\label{tab:arcscaling}
\small
\begin{tabular}{@{}rr|rrr|r@{}}
\toprule
 & & \multicolumn{3}{c|}{$64$-bit distances} & $128$-bit\\
$n$ & $k$ & \textsc{DP} & \textsc{O1} & \textsc{O2} & \textsc{O2}\\
\midrule
\tabinput{generated/arc_scaling}
\bottomrule
\end{tabular}
\end{table}





\chapter{Conclusion and Future Work}\label{ch:conclusion}

\begin{chapterabstract}
This chapter summarises the results of the project, states their
limitations, and lists directions for future work.
\end{chapterabstract}

\section{Conclusions}

The project studied $k$-dispersion for points in convex position, starting
from the algorithms of Singireddy, Basappa and Mitchell~\citep{sbm}. Its main
outcomes are the following.

\begin{itemize}[leftmargin=*,itemsep=2pt]
\item The bounded search tree of~\citet{sbm} is correct for $k\le4$
(\cref{thm:sbm4}). The proof combines a Delaunay ear argument, which makes the
first step safe for every $k$ (\cref{thm:firststep}), with a lemma on convex
quadrilaterals and a five-point lemma proved by angle chasing. Natural rules
for choosing the starting point fail for $k=5,6,7$ (\cref{prop:rootrules}).
\item The greedy properties behind the fast algorithms for $k=3$ hold for
triples and fail for every $k\ge4$, with explicit small counterexamples
(\cref{sec:greedy}).
\item Under the arc property a single greedy chain is exact
(\cref{thm:arc}); for unimodal point sets this gives an
$O(n\log^2n)$ expected-time algorithm (\cref{thm:arcopt}), compared with the
cubic bounds for general input.
\item A Delaunay certificate for $r$-far sets yields a short correctness proof of
a capped dynamic program that decides the problem in $O(n^3\log k)$ time and
solves it in $O(n^3\log k\log n)$ time (\cref{ch:exact}).
\item All algorithms were implemented with exact arithmetic and agreed with
exhaustive search on every test instance (\cref{ch:results}).
\end{itemize}

\section{Limitations}\label{sec:limitations}

\paragraph{The two-sided search tree.}
$(\mathrm{SBM}_k)$ remains unproved for $k\ge5$, so the bound
$O(2^kn^2\log^2n)$ of~\citet{sbm} is established here only for $k\le4$; the
published argument for general $k$ is a picture argument with three centres of
a Delaunay triangle that we have not been able to complete. Our method does not
extend directly: after a replacement the new point moves away from the
remaining witness points, so the attachment condition of \cref{lem:attach} is
not inherited, and \cref{prop:rootrules} shows that the root must be chosen
with more care than by distances or angles alone. Not even every vertex of a
solution is a good root for $k=4$ (the last instance in
\code{fpt/data/root\_rule\_instances.json}). We conjecture the following
strengthening of $(\mathrm{SBM}_k)$, which holds for $k\le4$ by
\cref{thm:sbm4}: \emph{for every $r$-far $k$-set $T$ there is a vertex of $T$
from which a search that only takes safe steps reaches $k$ points.} The
evidence is experimental: the conjecture held on more than $1.2\cdot10^7$
pairs $(T,r)$ with $k\le7$, including adversarially optimised instances, and
the two-sided search succeeded in all $532{,}833$ random tests and about
$3{,}000$ adversarial runs for $k\le10$. Until a proof is found,
\cref{alg:certified} uses the search only as a sound heuristic.

\paragraph{The arc property.}
The near-linear bound of \cref{thm:arcopt} requires a unimodal point set; for
points exactly on a circle a linear-time algorithm was already
known~\citep{akagi}. For general convex input the arc property must be checked
at every tested radius, which costs $O(n^2)$ per radius, and the adaptive
algorithm may report that it is not applicable. The arc property fails on thin
inputs: in our experiments it held at the optimal radius on every hull of a
random point cloud and on near-circular sets, but never on thin ellipses and
lenses, and applying the decision procedure without the hypothesis gave wrong
answers in $\arcUncheckedWrong\%$ of the cases on thin families. The question
whether contiguous far sets for the points of one solution suffice is open for
$k=4,5$ (\cref{prop:weaker} settles $k=6$).

\paragraph{The exact algorithm.}
The worst-case improvement over~\citet{tw26} is a factor $\log n/\log k$, which
vanishes when $k=n^{\Omega(1)}$, and the per-chord monotonicity behind it is
that of~\citet{tw26}; their algorithm also handles weighted disks of arbitrary
radii, which ours does not. In our experiments the cap made the decision
procedure $\capSpeedupKFour$ times faster for $k=4$ but only $1.1$--$1.4$ times
faster for larger $k$, and the measured growth was $n^{\expCappedFour}$ to
$n^{\expCappedSixteen}$ rather than $n^3$, because the number of binary-search
steps grows with $\log n$ and the table of $\Theta(n^2k)$ entries leaves the
cache. The certified algorithm is no faster than the exact one in the worst
case: when the heuristic bound is not tight, which happened only on thin
inputs, it was about $\certOverhead\%$ slower.

\paragraph{Experiments.}
All experiments use synthetic instances (hulls of random point clouds, points
on or near circles, thin ellipses and lenses) on a single machine. Integer
coordinates are bounded by $2^{29}$ so that the predicates fit into $128$ bits;
unimodal instances with $10^5$ or more points do not fit into this range and
were run with $128$-bit squared distances.



\section{Future work}

\begin{itemize}[leftmargin=*,itemsep=2pt]
\item Prove the claim of~\citet{sbm} for $k\ge5$, or find a counterexample.
The conjectured strengthening stated in \cref{sec:limitations}, that some vertex
of every $r$-far $k$-set is a good root for a search that only takes safe
steps, would suffice.
\item Close the gap between the arc property for all points and for the points
of one solution for $k=4,5$, and find further classes of point sets, beyond
unimodal ones, on which greedy is exact.
\item Remove the remaining $\log n$ factor of the optimisation, for example by
parametric search or by searching directly over the thresholds of the
dynamic program.
\item Extend the capped dynamic program to weighted dispersion and to the
maximum-weight independent set of disks considered by~\citet{tw26}.
\item Evaluate the algorithms on real data, such as candidate facility sites
along a coastline.
\end{itemize}


%%%%%%%%%%%%%%%%%%%%%%%%%%%%%%%%%%%%%%%%%%%%%%%%%%%%%%%%%%%%%%%%%%%%
\chapter*{Appendix}
\addcontentsline{toc}{chapter}{Appendix}
\renewcommand{\thesection}{A.\arabic{section}}
\setcounter{section}{0}
\section{Point sets without tight pairs}\label{app:instances}

\Cref{tab:instances} lists the instances behind \cref{thm:tightpair} for
$k=5,6,7$, together with two small instances that separate weaker properties
from $(\mathrm{TP}_k)$. All of them are contained in
\code{convex/data/counterexamples.json} and are verified by
\code{convex/experiments/check_counterexamples.py} by enumerating all
$k$-subsets in exact integer arithmetic.

\begin{table}[H]
\centering
\caption{Explicit instances (points in ccw order). In the first three, no
$r$-far $k$-set contains a tight pair. In the last two an $r$-far $k$-set
with a tight pair exists, but the stated rule fails.}
\label{tab:instances}
\small
\begin{tabular}{@{}rrrp{0.72\textwidth}@{}}
\toprule
$k$ & $n$ & $r^2$ & points\\
\midrule
5 & 10 & 1000000 & \raggedright $(-512,859)$, $(-995,-25)$, $(-1009,-52)$, $(-10,-1)$, $(2,0)$, $(1006,163)$, $(1367,222)$, $(554,833)$, $(518,858)$, $(-492,864)$\tabularnewline\addlinespace
6 & 12 & 418609 & \raggedright $(851,88)$, $(215,259)$, $(-175,352)$, $(-820,491)$, $(-1031,524)$, $(-687,-38)$, $(-649,-92)$, $(-13,-263)$, $(377,-356)$, $(1024,-495)$, $(1233,-528)$, $(889,34)$\tabularnewline\addlinespace
7 & 14 & 13155129 & \raggedright $(5936,1558)$, $(2254,1872)$, $(-36,2004)$, $(-3872,2121)$, $(-4826,2102)$, $(-7427,-517)$, $(-8045,-1226)$, $(-4359,-1442)$, $(-3984,-1457)$, $(763,-1433)$, $(3767,-1353)$, $(7528,-1170)$, $(8805,-1040)$, $(6069,1446)$\tabularnewline
\midrule
4 & 8 & 50 & \raggedright $(\mathrm{TT}_4)$ fails, $(\mathrm{ST}_4)$ holds: $(-11,0)$, $(-9,-1)$, $(1,-2)$, $(3,-2)$, $(10,0)$, $(8,1)$, $(1,2)$, $(-4,2)$\tabularnewline\addlinespace
6 & 10 & 26 & \raggedright $(\mathrm{AG}_6)$ fails, $(\mathrm{ST}_6)$ holds: $(-13,0)$, $(-9,-1)$, $(0,-2)$, $(2,-2)$, $(8,-1)$, $(13,0)$, $(11,1)$, $(1,2)$, $(-2,2)$, $(-8,1)$\tabularnewline
\bottomrule
\end{tabular}
\end{table}



\end{spacing}

\clearpage
\renewcommand{\bibname}{References}
\bibliographystyle{apalike}
\addcontentsline{toc}{chapter}{References}
\bibliography{mybib}

\end{document}
